We now assume that \(\varphi_{\lambda}\) does not tend to 0 at infinity, and we denote by \(\theta = 1 / \lambda\) . For every real number \(x\) , we denote by \(\langle x\rangle\) the unique real number in \([0,1[\) such that \(x - \langle x\rangle \in \mathbf{Z}\) .

\begin{enumerate}
\def\labelenumi{\alph{enumi})}
\setcounter{enumi}{7}
\item
  There exists a real number \(u \neq 0\) such that the series with general term \(\sin(u\theta^{n})^{2}\) converges; the series with general term \(\langle u\theta^{n}\rangle\) converges.
\item
  For every \(n \in N\) , let \(c_{n}\) be an integer such that \(|u\theta^{n} - c_{n}| \leqslant \frac{1}{2}\) . The formal series \(g(z) = \sum c_{n}z^{n}\) is a rational function. (Use the criterion from A, IV, p.~85, Exercise 1, and Hadamard's inequality, EVT, V, p.~37, Cor. 3, to show that the Hankel determinants that arise tend to 0; conclude by observing that these determinants are integers.)
\item
  There exist polynomials with integer coefficients \(f_1\) and \(f_2\) such that \(g = f_1 / f_2\) and \(f_2(0) = 1\) .
\item
  The number \(\theta\) is a Pisot number (“Salem’s theorem”). (Note that \(f_{2}(\lambda)=0\) and that the other roots of \(f_{2}\) have modulus \textgreater1.)
\end{enumerate}

\begin{enumerate}
\def\labelenumi{\arabic{enumi})}
\setcounter{enumi}{58}
\tightlist
\item
  Let \(n \geqslant 1\) be an integer and equip \(R^{n}\) with the Euclidean norm. We denote by \(\mu\) the Lebesgue measure on \(R^{n}\) .
\end{enumerate}

Let r \textgreater{} 0. A packing of spheres of radius r in \(R^{n}\) is a set P of disjoint subsets of \(R^{n}\) , such that every \(S \in P\) is a closed ball of radius r. We say that P is periodic if there exists a lattice \(\Lambda \subset R^{n}\) such that \(S + h \in P\) for all \(S \in P\) and every \(h \in \Lambda\) .

Let \(A _{P}\) be the union of the elements of P. It is a measurable set. We call the density of P the real number

\[
\Delta (\mathcal {P}) = \limsup _ {\mathrm{R} \to + \infty} \frac {\mu (\mathrm{A} _ {\mathcal {P}} \cap \mathrm{B} _ {\mathrm{R}})}{\mu (\mathrm{B} _ {\mathrm{R}})}
\]

where \(B_{R}\) is the closed ball of radius R in \(R^{n}\) . We denote by \(\Delta_{n} = \sup_{\mathcal{P}} \Delta(\mathcal{P})\) . a) We have

\[
\Delta_ {n} = \sup _ {\mathcal {P} \text {   periodic }} \Delta (\mathcal {P}).
\]

\begin{enumerate}
\def\labelenumi{\alph{enumi})}
\setcounter{enumi}{1}
\tightlist
\item
  Let P be a periodic sphere packing, invariant under the lattice \(\Lambda\) . Let \(C \subset R^{n}\) be a set of representatives of the centers of the balls \(S \in P\) modulo \(\Lambda\) . We have
\end{enumerate}

\[
\Delta (\mathcal {P}) = \mu (\mathrm{B} _ {1}) \frac {r ^ {n} \operatorname{Card} (\mathrm{C})}{\mathrm{V} (\Lambda)}
\]

where V(Λ) is the covolume of Λ.

\begin{enumerate}
\def\labelenumi{\alph{enumi})}
\setcounter{enumi}{2}
\tightlist
\item
  Let \(a > 0\) be a real number. Let \(f: \mathbf{R}^{n} \to \mathbf{R}\) be a Schwartz function such that
\end{enumerate}

\begin{enumerate}
\def\labelenumi{(\roman{enumi})}
\item
  We have \(f(x) \leqslant 0 \text{ if } \|x\| \geqslant 2a\) ;
\item
  The Fourier transform of \(f\) is real-valued and positive, and satisfies \(\widehat{f}(0)>0\) .
\end{enumerate}

Then we have

\[
\Delta_ {n} \leqslant \mu (\mathrm{B} _ {1}) a ^ {n} \frac {f (0)}{\widehat {f} (0)}.
\]

(It suffices to consider periodic sphere packings of radius \(a\) ; let \(\mathcal{P}\) be one of them, and let \(\Lambda\) and \(C \subset \mathbf{R}^n\) be as in the previous question; using the Poisson formula, prove that

\[
\sum_ {(c, d) \in \mathrm{C} ^ {2}} \sum_ {n \in \Lambda} f (d - c + n) = \frac {1}{\mathrm{V} (\Lambda)} \sum_ {h \in \Lambda^ {*}} \left| \sum_ {c \in \mathrm{C}} e ^ {2 i \pi c \cdot h} \right| ^ {2} \widehat {f} (h),
\]

where \(\Lambda^{*}\) is the associated lattice, and conclude.)

\begin{enumerate}
\def\labelenumi{\arabic{enumi})}
\setcounter{enumi}{59}
\tightlist
\item
  \begin{enumerate}
  \def\labelenumii{\alph{enumii})}
  \tightlist
  \item
    Let \(\mathbf{H} \subset \mathbf{C}\) be the open set of complex numbers \(z\) such that \(\mathcal{I}(z) > 0\) . The function
  \end{enumerate}
\end{enumerate}

\[
\Theta (z) = \sum_ {n \in \mathbf {Z}} e ^ {2 i \pi n ^ {2} z}
\]

is holomorphic on H, and it satisfies \(\Theta(z) = (z/i)^{-1/2}\Theta(-1/z)\) for \(z\in\mathrm{H}\) , where \(z\mapsto(z/i)^{1/2}\) is the unique holomorphic function on H that extends the function \(iy\mapsto\sqrt{y}\) . (Use Exercise 16, b).

\begin{enumerate}
\def\labelenumi{\alph{enumi})}
\setcounter{enumi}{1}
\tightlist
\item
  Let m and n be strictly positive integers. We have
\end{enumerate}

\[
\frac {1}{\sqrt {m}} \sum_ {k = 0} ^ {m - 1} e ^ {2 i \pi k ^ {2} n / m} = \frac {e ^ {i \pi / 4}}{\sqrt {2 n}} \sum_ {k = 0} ^ {2 n - 1} e ^ {- 2 i \pi k ^ {2} m / (2 n)}.
\]

(Apply the relation from the previous question with \(z = -2in / m + iy\) and let \(y \to 0\) .)

\begin{enumerate}
\def\labelenumi{\alph{enumi})}
\setcounter{enumi}{2}
\tightlist
\item
  For every integer \(m \geqslant 1\) and every \(n \in Z\) , set
\end{enumerate}

\[
\tau (n, m) = \sum_ {k = 0} ^ {m - 1} e ^ {2 i \pi k ^ {2} n / m}.
\]

For all strictly positive integers \(m_{1}\) and \(m_{2}\) we have

\[
\tau (1, m _ {1} m _ {2}) = \tau (m _ {2}, m _ {1}) \tau (m _ {1}, m _ {2}).
\]

\begin{enumerate}
\def\labelenumi{\alph{enumi})}
\setcounter{enumi}{3}
\item
  If m is odd, we have \(\tau(1,m)=i^{((m-1)/2)^{2}}\sqrt{m}\) . (Use a).
\item
  If p is an odd prime number, then \(\tau(n,p)=\ell_{p}(n)\tau(1,p)\) , where \(\ell_{p}\) is the unique character of order 2 of \((\mathbf{Z}/p\mathbf{Z})^{*}\) , extended to \(Z/pZ\) by setting \(\ell_{p}(0)=0\) .
\item
  Let \(p \neq q\) be two odd prime numbers. We have
\end{enumerate}

\[
\ell_ {p} (q) \ell_ {q} (p) = \frac {\tau (1 , p q)}{\tau (1 , p) \tau (1 , q)}.
\]

\begin{enumerate}
\def\labelenumi{\alph{enumi})}
\setcounter{enumi}{6}
\tightlist
\item
  Deduce from this a new proof of the law of quadratic reciprocity (cf. A, V, p. 156, Exercise 23): for all odd prime numbers \(p \neq q\) , one has
\end{enumerate}

\[
\ell_ {p} (q) \ell_ {q} (p) = (- 1) ^ {(p - 1) (q - 1) / 4}.
\]

\begin{enumerate}
\def\labelenumi{\arabic{enumi})}
\setcounter{enumi}{60}
\tightlist
\item
  Let \(r \geqslant 1\) be an integer. Let \(A = (a_{i,j})\) be a symmetric positive definite matrix of type \((r, r)\) with integer coefficients, and let \(N \geqslant 1\) be an integer such that the matrix \(A^* = NA^{-1}\) also has integer coefficients. We denote by Q (resp. \(\widetilde{Q}\) ) the integral quadratic form associated with A (resp. with \(A^{-1}\) ), that is, one has \(Q(x) = {}^t xAx\) for all \(x \in \mathbf{Z}^r\) . We denote by \((a, b) \mapsto \langle a|b \rangle_A\) the bilinear form associated with A, so that \(\langle a|b \rangle_A = {}^t aAb\) .
\end{enumerate}

We denote by H the set of complex numbers z such that \(\mathcal{I}(z)>0\) . This is an open subset of C. We denote by log: \(H\to C\) the restriction to H of the principal determination of the logarithm (FVR, III, p.~10). Let \(k\in N\) be an odd integer. For \(z\in H\) , we denote by \(z^{k}=\exp(k\log(z))\) .

\begin{enumerate}
\def\labelenumi{\alph{enumi})}
\tightlist
\item
  The group \(\mathrm{SL}_2(\mathbf{Z})\) acts transitively on \(\mathbf{H}\) by
\end{enumerate}

\[
\left( \begin{array}{c c} a & b \\ c & d \end{array} \right) \cdot z = \frac {a z + b}{c z + d}, \qquad \left( \begin{array}{c c} a & b \\ c & d \end{array} \right) \in \mathrm{SL} _ {2} (\mathbf {Z})
\]

  (cf. TA, I, p. 150, Exercise 5).

\begin{enumerate}
\def\labelenumi{\alph{enumi})}
\setcounter{enumi}{1}
\tightlist
\item
  For every complex number \(z \in \mathbf{H}\) and every \(x \in \mathbf{Z}^r\) , the series
\end{enumerate}

\[
\sum_{\substack{m\in \mathbf{Z}^{r}\\ m\equiv x\bmod \mathrm{N}}}\exp (i\pi \mathrm{Q}(m)z)
\]

converges absolutely and uniformly on compact sets.

We denote by \(\theta_{\mathrm{A}}(z;x)\) or \(\theta_{\mathrm{Q}}(z;x)\) the sum of this series (“theta function associated with Q”). It is a holomorphic function on \(\mathbf{H}\) .

\begin{enumerate}
\def\labelenumi{\alph{enumi})}
\setcounter{enumi}{2}
\tightlist
\item
  For every \(y \in \mathbf{Z}^r\) , one has
\end{enumerate}

\[
\sum_ {m \in {\bf Z} ^ {r}} \exp (i \pi Q (m + y) z) = \frac {1}{\sqrt {\det (A)}} \Bigl (\frac {i}{z} \Bigr) ^ {r / 2} \sum_ {m \in {\bf Z} ^ {r}} \exp \Bigl (- i \pi \frac {\widetilde {Q} (m)}{z} + m \cdot y \Bigr),
\]

(Apply Poisson's formula and Exercise 1.)

\begin{enumerate}
\def\labelenumi{\alph{enumi})}
\setcounter{enumi}{3}
\tightlist
\item
  We denote by H the subgroup of \((\mathbf{Z}/\mathrm{NZ})^{r}\) consisting of the elements \(\alpha\in(\mathbf{Z}/\mathrm{NZ})^{r}\) such that \(A\alpha=0\) in Z/NZ. The group H has cardinality det(A). Every character of H is of the form
\end{enumerate}

\[
\psi_ {\beta} (\alpha) = \exp \Bigl (2 i \pi \frac {\langle \alpha | \beta \rangle_ {\mathrm{A}}}{\mathrm{N} ^ {2}} \Bigr),
\]

where \(\beta\in\mathrm{H}\) .

\begin{enumerate}
\def\labelenumi{\alph{enumi})}
\setcounter{enumi}{4}
\item
  For every \(\alpha \in \mathrm{H}\) and \(z\in \mathbf{H}\) , one has \(\theta_{\mathrm{Q}}(z + 2;\alpha) = \exp \bigl (\frac{2i\pi\mathrm{Q}(\alpha)}{\mathrm{N}^2}\bigr)\theta_{\mathrm{Q}}(z;\alpha)\) .
\item
  For every \(\alpha \in \mathrm{H}\) and \(z\in \mathbf{H}\) , one has
\end{enumerate}

\[
\theta_ {\mathrm{Q}} \Bigl (- \frac {1}{z}; \alpha \Bigr) = \frac {1}{\sqrt {\det (\mathrm{A})}} (- i z) ^ {r / 2} \sum_ {\beta \in \mathrm{H}} \psi_ {\beta} (\alpha) \theta_ {\mathrm{Q}} (z; \beta).
\]

\(g)\) Let \(g = \begin{pmatrix} a & b \\ c & d \end{pmatrix} \in \mathrm{SL}_2(\mathbf{Z})\) such that \(c \neq 0\) , \(d > 0\) and \(b \equiv c \equiv 0 \mod 2\) . For \(\alpha \in \mathrm{H}\) and \(z \in \mathbf{H}\) , one has

\[
\theta_ {\mathrm{Q}} (g \cdot z; \alpha) = \frac {1}{d ^ {r / 2} \det (\mathrm{A})} (c z + d) ^ {r / 2} \sum_ {\beta \in \mathrm{H}} \Phi (\alpha , \beta) \theta_ {\mathrm{Q}} (z; \beta),
\]

where

\[
\Phi (\alpha ,\beta) = \sum_{\gamma \in \mathrm{H}}\psi_{\gamma}(\alpha)\Bigl (\sum_{\substack{\delta \text{mod.} d\mathrm{N}\\ \delta \equiv \alpha \text{mod.}\mathrm{N}}}\exp \Bigl (\frac{i\pi}{d\mathrm{N}^{2}} (b\mathrm{Q}(\delta) + 2\langle \gamma |\delta \rangle_{\mathrm{A}} - c\mathrm{Q}(\gamma))\Bigr)\Bigr)
\]

(First consider \(\theta_{\mathrm{Q}}(\widetilde{g} \cdot z; \alpha)\) where \(\widetilde{g} = g\begin{pmatrix} 0 & -1 \\ 1 & 0 \end{pmatrix}\) . Note the formula \(d\widetilde{g} \cdot z = b - (dz - c)^{-1}\) and apply the preceding question; note that \(\Phi(\alpha, \beta)\) depends only on \(\beta\) modulo N).

\begin{enumerate}
\def\labelenumi{\alph{enumi})}
\setcounter{enumi}{7}
\tightlist
\item
  Suppose moreover that \(c \equiv 0 \mod. 2\mathrm{N}\) . Then
\end{enumerate}

\[
\theta_ {\mathrm{Q}} (g \cdot z; x) = \frac {\tau}{d ^ {r / 2}} (c z + d) ^ {k} \theta_ {\mathrm{Q}} (z; a \alpha),
\]

where

\[
\tau = \sum_{\substack{\delta \text{mod.} d\mathrm{N}\\ \delta \equiv x\text{mod.}\mathrm{N}}}\exp \Bigl (\frac{i\pi b\mathrm{Q}(\delta)}{d\mathrm{N}^{2}}\Bigr).
\]

\begin{enumerate}
\def\labelenumi{\roman{enumi})}
\tightlist
\item
  Suppose moreover that \(d\) is odd. We have \(\tau = \exp \left(\frac{i\pi abQ(\alpha)}{N^2}\right)\tau_Q(c;d)\) , where
\end{enumerate}

\[
\tau_ {\mathrm{Q}} (c, d) = \sum_ {x \in (\mathbf {Z} / d \mathbf {Z}) ^ {r}} e \left(- \frac {i \pi c \mathrm{Q} (x)}{d}\right)
\]

(“Gauss sum.”)

\begin{enumerate}
\def\labelenumi{\alph{enumi})}
\setcounter{enumi}{9}
\tightlist
\item
  Let us write \(d = p_{1} \cdots p_{k}\) , where \(p_{1}, \ldots, p_{k}\) are prime numbers. We denote
\end{enumerate}

\[
\left(\frac {n}{d}\right) = \prod_ {i = 1} ^ {k} \ell_ {p _ {i}} (n)
\]

where, for every odd prime number \(p\) , we denote by \(\ell_p\) the unique character of order 2 of \((\mathbf{Z}/p\mathbf{Z})^*\) , extended to \(\mathbf{Z}/p\mathbf{Z}\) by setting \(\ell_p(0) = 0\) (“Jacobi symbol”). If \(d\) and \(2c\det(\mathrm{A})\) are coprime, then

\[
\tau_ {\mathrm{Q}} (c, d) = \Big (\frac {\det (\mathrm{A})}{d} \Big) \Big (\overline {{\varepsilon}} _ {d} \Big (\frac {2 c}{d} \Big) \sqrt {d} \Big) ^ {r},
\]

where

\[
\varepsilon_ {d} = \left\{ \begin{array}{l l} 1 & \text { if } d \equiv 1 \bmod . 4 \\ i & \text { if } d \equiv - 1 \bmod . 4. \end{array} \right.
\]

(Diagonalize the quadratic form and use the preceding exercise.)

\(k)\) If \(g\) is congruent to the identity modulo 4N, then

\[
\theta_ {\mathrm{Q}} (g \cdot z; \alpha) = \left(\frac {2 c}{d}\right) ^ {r} (c z + d) ^ {k} \theta_ {\mathrm{Q}} (z; \alpha).
\]

\(^{*}\) This means that the map \(z \mapsto \theta_{q}(z; \alpha)\) is a holomorphic modular form of weight r/2 and level 4N. (FM, in preparation.) \(^{*}\)

\begin{enumerate}
\def\labelenumi{\arabic{enumi})}
\setcounter{enumi}{61}
\tightlist
\item
  Let us identify \(\mathbf{T}\) with the interval \([-1/2, 1/2]\) . Let \(f\) be the function on \(\mathbf{R}/\mathbf{Z}\) defined by \(f(x) = s(x)\) (sign function). Let \(s_n\) be the symmetric partial sum
\end{enumerate}

\[
s _ {n} (x) = \sum_ {h = - n} ^ {n} \widehat {f} (h) e ^ {2 i \pi h x}
\]

of the Fourier series of f.

\begin{enumerate}
\def\labelenumi{\alph{enumi})}
\item
  For every \(x \neq 0\) in \(\mathbf{T}\) , one has \(s_n(x) \to f(x)\) when \(n \to +\infty\) .
\item
  Let \(n \geqslant 1\) and \(0 < x \leqslant 1/2\) . We have
\end{enumerate}

\[
s _ {2 n} (x) = \frac {\sin (2 \pi n x)}{2 \sin (\pi x)}.
\]

\begin{enumerate}
\def\labelenumi{\alph{enumi})}
\setcounter{enumi}{2}
\tightlist
\item
  We have
\end{enumerate}

\[
\lim _ {n \rightarrow + \infty} s _ {2 n} \left(\frac {1}{4 n}\right) = 2 \int_ {0} ^ {1 / 2} \frac {\sin (2 \pi x)}{\pi x} d x,
\]

and this limit belongs to the interval \(]1,2[\) (“Gibbs phenomenon”).

\begin{enumerate}
\def\labelenumi{\arabic{enumi})}
\setcounter{enumi}{62}
\tightlist
\item
  Let G be the subgroup of \(T^{N}\) consisting of the elements \(x = (x_{n})\) such that \(2x_{n} = 0\) for all \(n \in N\) except for at most finitely many.
\end{enumerate}

\begin{enumerate}
\def\labelenumi{\alph{enumi})}
\item
  For every finite set \(\Lambda \subset \mathbf{N}\) , let \(\mathrm{U}_{\Lambda} \subset \mathrm{G}\) be the set of \((x_{n})\) such that \(x_{n} = 1\) for \(n \in \Lambda\) and \(2x_{n} = 0\) for \(n \notin \Lambda\) . There exists a topological group structure on G such that the family of sets \(U_{\Lambda}\) is a base of open neighborhoods of e.
\item
  The topological group G is locally compact.
\item
  The group G is not divisible.
\end{enumerate}

\(d)\) The dual group \(\widehat{\mathrm{G}}\) is torsion-free.

\begin{enumerate}
\def\labelenumi{\arabic{enumi})}
\setcounter{enumi}{63}
\tightlist
\item
  Let \(\mu\) be a measure on \(\mathbf{T}\) .
\end{enumerate}

\begin{enumerate}
\def\labelenumi{\alph{enumi})}
\tightlist
\item
  For every \(x \in \mathbf{T}\) , one has
\end{enumerate}

\[
\mu (\{x \}) = \lim _ {\mathrm{H} \rightarrow + \infty} \frac {1}{2 \mathrm{H} + 1} \sum_ {| h | \leqslant \mathrm{H}} \widehat {\mu} (h) e ^ {2 i \pi h x}.
\]

\begin{enumerate}
\def\labelenumi{\alph{enumi})}
\setcounter{enumi}{1}
\tightlist
\item
  Let us write \(\mu = \mu_d + \mu_a\) , where \(\mu_a\) is an atomic measure and \(\mu_d\) is a diffuse measure (INT, V, §5, no. 10, Prop. 15). Let S be the support of \(\mu_a\) . We have
\end{enumerate}

\[
\lim _ {\mathrm{H} \to + \infty} \frac {1}{2 \mathrm{H} + 1} \sum_ {| h | \leqslant \mathrm{H}} | \widehat {\mu} (h) | ^ {2} = \sum_ {x \in \mathrm{S}} | \mu (\{s \}) | ^ {2}
\]

(consider \((\mu * \widetilde{\mu})(\{0\})\) ).

\begin{enumerate}
\def\labelenumi{\alph{enumi})}
\setcounter{enumi}{2}
\tightlist
\item
  Suppose that \(\mu\) is a probability measure on \(\mathbf{T}\) . The following conditions are equivalent:
\end{enumerate}

\begin{enumerate}
\def\labelenumi{(\roman{enumi})}
\item
  There exists \(x \in \mathbf{T}\) such that \(\mu = \varepsilon_x\) ;
\item
  We have \(\lim_{|h|\to+\infty}|\widehat{\mu}(h)|^{2}=1\) ;
\item
  We have
\end{enumerate}

\[
\lim _ {\mathrm{H} \rightarrow + \infty} \frac {1}{2 \mathrm{H} + 1} \sum_ {| h | \leqslant \mathrm{H}} | \widehat {\mu} (h) | ^ {2} = 1.
\]

\begin{enumerate}
\def\labelenumi{\arabic{enumi})}
\setcounter{enumi}{64}
\tightlist
\item
  Assume that G is compact and equipped with its normalized Haar measure \(\mu\) . Let \(f: G \to G\) be a continuous and surjective homomorphism.
\end{enumerate}

\begin{enumerate}
\def\labelenumi{\alph{enumi})}
\item
  We have \(f(\mu) = \mu\) .
\item
  The map \(f\) is \(\mu\) -ergodic (cf. exercise 16 of I, p. 190) if and only if there exists no character \(\chi \neq e_{\widehat{\mathrm{G}}}\) of G and integer \(n \geqslant 1\) such that \(\chi \circ f^n = \chi\) .
\item
  Let \(d \geqslant 1\) be an integer and let a be a matrix of type \((d, d)\) with integer coefficients and determinant 1. The map \(f_{a} \colon T^{d} \to T^{d}\) defined by \(f_{a}(x) = ax\) for all \(x \in T^{d}\) is ergodic with respect to the normalized Haar measure on \(T^{d}\) if and only if the spectrum of the matrix a in \(\mathcal{L}(\mathbf{C}^{n})\) contains no root of unity.
\end{enumerate}

¶ 66) We denote \(e\colon \mathbf{T}\to \mathbf{C}^{*}\) the character defined by passing to the quotient by the character \(x\mapsto e^{2i\pi x}\) of \(\mathbf{R}\) .

Let \(d \geqslant 2\) be an integer. Let \(\alpha \in R\) and let \(g \in R[X]\) be a polynomial whose degree is \textless{} d. Let a and q be coprime integers, with \(q \neq 0\) , such that

\[
\left| \alpha - \frac {a}{q} \right| \leqslant \frac {1}{q ^ {2}}.
\]

\begin{enumerate}
\def\labelenumi{\alph{enumi})}
\tightlist
\item
  Suppose that \(\alpha \notin \mathbf{Z}\) . Let N and M be integers with \(\mathrm{M} \geqslant 0\) . We have
\end{enumerate}

\[
\Big | \sum_ {n = N} ^ {N + M} e ^ {2 i \pi \alpha n} \Big | \leqslant \frac {1}{| \sin (\pi \alpha) |}.
\]

\begin{enumerate}
\def\labelenumi{\alph{enumi})}
\setcounter{enumi}{1}
\tightlist
\item
  We denote \(D = 2^{d-1}\) . For every \(\varepsilon > 0\) , there exists a real number \(C(\varepsilon, d) \geqslant 0\) such that for every integer \(P \geqslant 1\) , one has
\end{enumerate}

\[
\left| \sum_ {n = 1} ^ {\mathrm{P}} e ^ {2 i \pi (\alpha n ^ {d} + g (n))} \right| \leqslant \mathrm{C} (\varepsilon , d) \mathrm{P} ^ {1 + \varepsilon} \left(\mathrm{P} ^ {- 1 / \mathrm{D}} + q ^ {- 1 / \mathrm{D}} + \left(\frac {\mathrm{P} ^ {d}}{q}\right) ^ {- 1 / \mathrm{D}}\right).
\]

(Proceed by induction on d, reducing the square of the left-hand side of the inequality to a family of polynomials of degree \(\leqslant d - 1\) ; when d = 2, apply part a). One may prove and use the bound

\[
d (n) = \mathrm{O} (n ^ {\varepsilon})
\]

when \(n \to +\infty\) , where \(n \geqslant 1\) and \(d(n)\) is the number of divisors of n in N, valid for every real number \(\varepsilon > 0\) .

\begin{enumerate}
\def\labelenumi{\alph{enumi})}
\setcounter{enumi}{2}
\tightlist
\item
  Deduce a new proof of the assertion of part c) of Exercise 35.
\end{enumerate}

¶ 67) Let \(k \geqslant 2\) and \(\ell > k\) be natural numbers. For every integer \(N \geqslant 1\) , we denote by \(r_{k,\ell}(N)\) the number of families \((n_1, \ldots, n_\ell) \in \mathbf{N}^\ell\) such that

\[
n _ {1} ^ {k} + \dots + n _ {\ell} ^ {k} = \mathrm{N}.
\]

The goal of this exercise is to show that, if \(\ell\) is sufficiently large as a function of \(k\) , then \(r_{k,\ell}(\mathrm{N}) \geqslant 1\) (“Waring's problem”) for every \(\mathrm{N} \geqslant 1\) .

We denote by \(e\colon\mathbf{T}\to\mathbf{C}^{*}\) the character defined by passing to the quotient by the character \(x\mapsto e^{2i\pi x}\) of \(\mathbf{R}\) .

Let P ≥ 1 be an integer. We denote by S \(_{P}\) the function on T such that

\[
\mathrm{S} _ {\mathrm{P}} (x) = \sum_ {n = 1} ^ {\mathrm{P}} e (x n ^ {k}).
\]

We denote by \(dx\) the normalized Haar measure on \(\mathbf{T}\) .

\begin{enumerate}
\def\labelenumi{\alph{enumi})}
\tightlist
\item
  If \(\mathrm{P} > \mathrm{N}^{1 / k}\) , then we have
\end{enumerate}

\[
r _ {k, \ell} (\mathrm{N}) = \int_ {\mathbf {T}} \mathrm{S} _ {\mathrm{P}} (x) ^ {\ell} e (- \mathrm{N} x) d x.
\]

\begin{enumerate}
\def\labelenumi{\alph{enumi})}
\setcounter{enumi}{1}
\tightlist
\item
  Let \(\delta\) be a real number such that \(0 < \delta < 1/3\) and let \(\mathrm{P} = [\mathrm{N}^{1/k}]\) . We assume that N is sufficiently large, so that \(\mathrm{P} \geqslant 1\) . Let \(a\) and \(q\) be pairwise coprime integers such that \(1 \leqslant q \leqslant \mathrm{P}^{\delta}\) . We denote by \(\mathrm{M}_{a,q}\) the image in \(\mathbf{T}\) of the open interval
\end{enumerate}

\[
\left. \right] \frac {a}{q} - \frac {1}{\mathrm{P} ^ {k - \delta}}, \frac {a}{q} + \frac {1}{\mathrm{P} ^ {k - \delta}} \Big [.
\]

For \(x\in \mathrm{M}_{a,q}\) , one has

\[
\sum_ {n = 1} ^ {\mathrm{P}} e ^ {2 i \pi x n ^ {k}} = \frac {1}{q} \mathrm{S} (a; q) \mathrm{I} \Bigl (x - \frac {a}{q} \Bigr) + \mathrm{O} (\mathrm{P} ^ {2 \delta})
\]

as N → +∞, where

\[
\mathrm{S} (a; q) = \sum_ {x = 1} ^ {q} e \left(\frac {a x ^ {k}}{q}\right), \quad \mathrm{I} (y) = \int_ {0} ^ {\mathrm{P}} e \left(y t ^ {k}\right) d t.
\]

\begin{enumerate}
\def\labelenumi{\alph{enumi})}
\setcounter{enumi}{2}
\tightlist
\item
  We denote by M the union of the sets \(M_{a,q}\), where q ranges over the integers such that \(1 \leqslant q \leqslant P^{\delta}\) and a ranges over the integers coprime to q such that \(1 \leqslant a < q\) . If \(\delta < 1/5\) , then there exists a real number \(\delta' > 0\) such that
\end{enumerate}

\[
\int_ {\mathrm{M}} \mathrm{S} _ {\mathrm{P}} (x) ^ {\ell} e (- \mathrm{N} x) d x = \mathrm{P} ^ {\ell - k} \mathfrak {S} (\mathrm{P} ^ {\delta}, \mathrm{N}) \mathrm{J} (\mathrm{P} ^ {\delta}) + \mathrm{O} (\mathrm{P} ^ {\ell - k - \delta^ {\prime}})
\]

as N → +∞, where we set, for Q ≥ 1

\[
\mathfrak{S}(\mathrm{Q},\mathrm{N}) = \sum_{1\leqslant q\leqslant \mathrm{Q}}\sum_{\substack{a = 1\\ (a,q) = 1}}^{q}\Big(\frac{1}{q}\mathrm{S}(a;q)\Big)^{\ell}e\Big(-\frac{aN}{q}\Big)
\]

\[
\mathrm{J(Q)} = \int_ {- \mathrm{Q}} ^ {\mathrm{Q}} \Bigl (\int_ {0} ^ {1} e (y t ^ {k}) d t \Bigr) ^ {\ell} e (- y) d y.
\]

\begin{enumerate}
\def\labelenumi{\alph{enumi})}
\setcounter{enumi}{3}
\tightlist
\item
  We have
\end{enumerate}

\[
\mathrm{J} (\mathrm{Q}) = \frac {\Gamma (1 + 1 / k) ^ {\ell}}{\Gamma (\ell / k)} + \mathrm{O} (\mathrm{Q} ^ {- (\ell / k - 1)})
\]

when Q → +∞. (First show that

\[
\mathrm{J} (\mathrm{Q}) = \frac {1}{k ^ {\ell}} \int_ {\mathrm{X} _ {\ell}} (y _ {1} \dots y _ {\ell - 1} (1 - y _ {1} - \dots - y _ {\ell - 1})) ^ {- 1 + 1 / k} d y + \mathrm{O} (\mathrm{Q} ^ {- (\ell / k - 1)})
\]

where

\[
\mathrm{X} _ {\ell} = \left\{\left(y _ {1}, \dots , y _ {\ell - 1}\right) \in \left(\mathbf {R} _ {+} ^ {*}\right) ^ {\ell - 1} \mid y _ {1} + \dots + y _ {\ell - 1} <   1 \right\},
\]

then evaluate this integral; cf.~FVR VII, p.~8, prop. 4.)

\begin{enumerate}
\def\labelenumi{\alph{enumi})}
\setcounter{enumi}{4}
\tightlist
\item
  Suppose that \(\ell >\delta^{-1}k2^{k}\) . There exists \(\delta^{\prime} > 0\) such that
\end{enumerate}

\[
\int_ {\mathbf {T - M}} | \mathrm{S} _ {\mathrm{P}} (x) | ^ {\ell} e (- \mathrm{N} x) d x = \mathrm{O} (\mathrm{P} ^ {\ell - k - \delta^ {\prime}})
\]

when \(N \rightarrow +\infty\) . (If N is sufficiently large and if \(x \in R\) has image in T-M, there exist coprime integers a and q such that

\[
\mathrm{P} ^ {\delta} \leqslant q \leqslant \mathrm{P} ^ {k - \delta}, \quad \left| x - \frac {a}{q} \right| \leqslant \frac {1}{q ^ {2}}.
\]

Then use Exercise 66.)

\begin{enumerate}
\def\labelenumi{\alph{enumi})}
\setcounter{enumi}{5}
\tightlist
\item
  For every prime number p and every integer \(m \geqslant 1\) , set
\end{enumerate}

\[
\mathrm{A}(p^{m}) = \sum_{\substack{a = 1\\ (a,q) = 1}}^{p^{m}}\Big(\frac{1}{p^{m}}\mathrm{S}(a;p^{m})\Big)^{\ell}e\Big(-\frac{a\mathrm{N}}{p^{m}}\Big)
\]

and

\[
\mathrm{N} (p ^ {m}) = \operatorname{Card} \left(\left\{\left(x _ {1}, \dots , x _ {\ell}\right) \in (\mathbf {Z} / p ^ {m} \mathbf {Z}) ^ {\ell} \mid x _ {1} ^ {k} + \dots + x _ {\ell} ^ {k} = \mathrm{N mod.} p ^ {m} \right\}\right).
\]

When \(\ell > 2^k\) , one has

\[
1 + \sum_ {m = 1} ^ {+ \infty} \mathrm{A} (p ^ {m}) = \lim _ {m \to + \infty} \frac {\mathrm{N} (p ^ {m})}{p ^ {m (\ell - 1)}}.
\]

\begin{enumerate}
\def\labelenumi{\alph{enumi})}
\setcounter{enumi}{6}
\tightlist
\item
  When \(\ell > 2^k\) , one has \(\mathfrak{S}(\mathrm{Q}, \mathrm{N}) = \mathfrak{S}(\mathrm{N}) + o(1)\) when \(\mathrm{Q} \to +\infty\) , where
\end{enumerate}

\[
\mathfrak {S} (\mathrm{N}) = \prod_ {p} \Bigl (1 + \sum_ {m \geqslant 1} \mathrm{A} (p ^ {m}) \Bigr),
\]

the product over the prime numbers being absolutely convergent.

\begin{enumerate}
\def\labelenumi{\alph{enumi})}
\setcounter{enumi}{7}
\tightlist
\item
  Let p be a prime number. Let \(\nu \geqslant 0\) be the integer such that \(k = p^{\nu}r\) where p does not divide r. Denote by \(\gamma\) the integer
\end{enumerate}

\[
\gamma = \left\{ \begin{array}{l l} \nu + 1 & \text { if } p \geqslant 3 \\ \nu + 2 & \text { otherwise. } \end{array} \right.
\]

If there exist integers \((n_1, \ldots, n_\ell)\) such that \(p\) does not divide all the \(n_i\) and such that

\[
x _ {1} ^ {k} + \dots + x _ {\ell} ^ {k} \equiv 0 \mathrm{mod.} p ^ {\gamma},
\]

then

\[
\lim _ {m \to + \infty} \frac {\mathrm{N} (p ^ {m})}{p ^ {m (\ell - 1)}} > 0.
\]

\begin{enumerate}
\def\labelenumi{\roman{enumi})}
\tightlist
\item
  There exists \(\ell > k\) such that, if N is sufficiently large, then \(\mathfrak{S}(N) > 0\) . (Use the previous question and part h) of Exercise 33.)
\end{enumerate}

\begin{enumerate}
\def\labelenumi{\alph{enumi})}
\setcounter{enumi}{9}
\tightlist
\item
  There exists an integer \(\ell > k\) such that \(r_{k,\ell}(\mathrm{N}) \geqslant 1\) for every integer \(\mathrm{N} \geqslant 1\) .
\end{enumerate}

\begin{enumerate}
\def\labelenumi{\arabic{enumi})}
\setcounter{enumi}{67}
\tightlist
\item
  We write the group law of G additively. Let \(\Gamma\) be a discrete subgroup of G such that \(G / \Gamma\) is compact. Let \(\Lambda = \Gamma^{\perp}\) be the orthogonal complement of \(\Gamma\) in \(\widehat{G}\) ; since \(\widehat{G / \Gamma}\) is identified with \(\Lambda\) , which is a discrete subgroup of \(\widehat{G}\) , and \(\widehat{G} / \Lambda\) is compact (since its dual identifies with \(\Gamma\) ).
\end{enumerate}

We assume that the Haar measure on G is chosen so that the quotient of the dual measure \(d\chi\) on \(\widehat{G}\) by the counting measure on \(\Lambda\) is the normalized Haar measure on \(\widehat{G}/\Lambda\) (INT, VII, § 2, no. 2, def. 1).

\begin{enumerate}
\def\labelenumi{\alph{enumi})}
\item
  There exists a measurable subset \(\Omega\) of \(\widehat{\mathbf{G}}\) which meets each class of \(\mathbf{G}\) modulo \(\Lambda\) at one point and only one; the measure of \(\Omega\) is equal to 1.
\item
  For \(x\in \mathbf{G}\) , set
\end{enumerate}

\[
\varphi (x) = \int_ {\Omega} \chi (x) d \chi .
\]

The function \(\varphi\) is continuous; it satisfies \(\varphi(0) = 1\) and \(\varphi(\gamma) = 0\) for all \(\gamma \in \Gamma - \{0\}\) .

\begin{enumerate}
\def\labelenumi{\alph{enumi})}
\setcounter{enumi}{2}
\tightlist
\item
  Prove that \(\varphi\) belongs to \(\mathrm{L}^2 (\mathrm{G})\) and that \(\| \varphi \| _2 = 1\) . For every \(\gamma \in \Gamma -\{e\}\) , one has
\end{enumerate}

\[
\int_ {\mathrm{G}} \varphi (x) \overline {{\varphi (\gamma - x)}} d x = 0.
\]

\begin{enumerate}
\def\labelenumi{\alph{enumi})}
\setcounter{enumi}{3}
\tightlist
\item
  Let \(f \in \mathrm{L}^{2}(\mathrm{G})\) such that \(\mathcal{F}(f)\) is zero almost everywhere outside of \(\Omega\) . Prove that
\end{enumerate}

\[
f = \sum_ {\gamma \in \Gamma} f (\gamma) \boldsymbol {\delta} _ {\gamma} (\varphi)
\]

in \(\mathrm{L}^2 (\mathrm{G})\) , where \(\delta_{\gamma}(\varphi)(x) = \varphi (x - \gamma)\) . Moreover, we have

\[
\int_ {\mathrm{G}} | f (x) | ^ {2} d x = \sum_ {\gamma \in \Gamma} | f (\gamma) | ^ {2}.
\]

\begin{enumerate}
\def\labelenumi{\alph{enumi})}
\setcounter{enumi}{4}
\tightlist
\item
  There exists a unique continuous function on G which coincides almost everywhere with f. If f is continuous, then
\end{enumerate}

\[
f = \sum_ {\gamma \in \Gamma} f (\gamma) \boldsymbol {\delta} _ {\gamma} (\varphi)
\]

in \(\mathcal{C}(\mathrm{G})\) .

\begin{enumerate}
\def\labelenumi{\alph{enumi})}
\setcounter{enumi}{5}
\tightlist
\item
  Specialize to the case where G = R and \(\Lambda = hZ\) for some h \textgreater{} 0 by taking \(\Omega = [-h, h[\) .
\end{enumerate}

\exercisesection{Classification}

\begin{enumerate}
\def\labelenumi{\arabic{enumi})}
\tightlist
\item
  \begin{enumerate}
  \def\labelenumii{\alph{enumii})}
  \tightlist
  \item
    For G to have a countable base, it is necessary and sufficient that \(\widehat{G}\) have a countable base. (If G has a countable base, then \(L^{1}(G)\) has a countable base, hence \(\mathsf{X}(\mathsf{L}^{1}(\mathsf{G}))\) has a countable base.)
  \end{enumerate}
\end{enumerate}

\begin{enumerate}
\def\labelenumi{\alph{enumi})}
\setcounter{enumi}{1}
\item
  For G to be countable at infinity, it is necessary and sufficient that \(\widehat{\mathrm{G}}\) admit an open subgroup with a countable base, or equivalently that \(\widehat{\mathrm{G}}\) be metrizable. (Use Prop. 3 of II, p.~248.)
\item
  Suppose G is compact. For G to be metrizable, it is necessary and sufficient that \(\widehat{G}\) be countable, or equivalently that G be isomorphic to a closed subgroup of \(T^{N}\) . (A countable discrete commutative group is a quotient of \(\mathbf{Z}^{(N)}\) .)
\end{enumerate}

\begin{enumerate}
\def\labelenumi{\arabic{enumi})}
\setcounter{enumi}{1}
\tightlist
\item
  \begin{enumerate}
  \def\labelenumii{\alph{enumii})}
  \tightlist
  \item
    Every infinite commutative group \(\Gamma\) admits an infinite countable quotient group. (Embed an infinite countable subgroup \(\Delta\) of \(\Gamma\) in a countable divisible group \(\Delta'\) by applying A, II, p.~185, Exercise 14; then extend to \(\Gamma\) the identity morphism of \(\Delta\) into \(\Delta'\) .)
  \end{enumerate}
\end{enumerate}

\begin{enumerate}
\def\labelenumi{\alph{enumi})}
\setcounter{enumi}{1}
\tightlist
\item
  Every infinite compact commutative group contains an infinite metrizable closed subgroup. (Use a).)
\end{enumerate}

\begin{enumerate}
\def\labelenumi{\arabic{enumi})}
\setcounter{enumi}{2}
\tightlist
\item
  \begin{enumerate}
  \def\labelenumii{\alph{enumii})}
  \tightlist
  \item
    Let K be a compact group and \(\mathrm{K}_{(n)}\) the set of elements of order n of K. If \(K \neq E_{n}\) for every n, then the set of elements of infinite order of K is dense in K. (If \(\mathrm{K}_{(n)}\) has nonempty interior, we have \(\mathrm{K}_{(m)} = \mathrm{K}\) for some integer \(m \geqslant n\) .)
  \end{enumerate}
\end{enumerate}

\begin{enumerate}
\def\labelenumi{\alph{enumi})}
\setcounter{enumi}{1}
\item
  Let \(\Gamma\) be a discrete commutative group whose elements are not of bounded order. Then \(\Gamma\) admits a countable quotient with the same property. (Reason as in part a) of Exercise 2.)
\item
  If every neighborhood of e in G contains an element of infinite order, then G possesses a metrizable closed subgroup with the same property. (Reduce to the case where G is compact, and use a) and b).) Otherwise, if G is not discrete, there exists an integer q \textgreater{} 1 such that G contains a closed subgroup isomorphic to \((\mathbf{Z}/q\mathbf{Z})^{\mathbf{N}}\) . (Use A, VII, p.~54, exerc. 4, c).)
\end{enumerate}

\begin{enumerate}
\def\labelenumi{\arabic{enumi})}
\setcounter{enumi}{3}
\item
  If G is compact and if \(\widehat{G}\) has an element of infinite order, then there exists a nontrivial continuous homomorphism from R to G. (Since R is divisible, there exists a nontrivial homomorphism from \(\widehat{G}\) into R.)
\item
  Let p be a prime number. The group \(Q_{p}\) is not the product of a compact group and a discrete group. (Every compact subgroup of \(Q_{p}\) is of the form \(p^{n}Z_{p}\) for some integer \(n \in Z\) .)
\item
  If G is a torsion group, then G and \(\widehat{G}\) are totally disconnected. (\(\widehat{G}\) is totally disconnected by Cor. 2 of II, p.~250; to prove that G is totally disconnected, reduce to the case where G is compact using Prop. 3 of II, p.~248, and then prove that \(\widehat{G}\) is a torsion group.)
\item
  An element \(x\) of \(G\) is said to be of infinite height if, for every \(n \in \mathbf{Z}\) , there exists \(y \in G\) such that \(y^n = x\) . Prove that, if \(G\) is compact, the set of elements of infinite height is the neutral component of \(G\) .
\end{enumerate}

¶ 8) The group G is locally connected if and only if G is isomorphic to a group \(\mathbf{R}^{n}\times\mathrm{E}\times\widehat{\mathrm{D}}\) where E and D are discrete and every finite-rank subgroup of D is free.

\begin{enumerate}
\def\labelenumi{\arabic{enumi})}
\setcounter{enumi}{8}
\tightlist
\item
  \begin{enumerate}
  \def\labelenumii{\alph{enumii})}
  \tightlist
  \item
    Let \(\boldsymbol{a}=(a_{n})_{n\geqslant0}\) be a sequence of integers \textgreater{} 1. We equip Z with the topological ring structure for which the subsets \(Za_{0}a_{1}\cdots a_{n}\) , where \(n\geqslant0\) , form a fundamental system of neighborhoods of 0. We denote by \(\Delta_{a}\) the completion of this ring. It is compact, metrizable, and totally disconnected. If p is a prime number, and if \(a_{i}=p\) for all i, we have \(\Delta_{a}=Z_{p}\) .
  \end{enumerate}
\end{enumerate}

\begin{enumerate}
\def\labelenumi{\alph{enumi})}
\setcounter{enumi}{1}
\item
  We denote by \(\mathbf{Z}(\boldsymbol{a}^{\infty})\) the multiplicative group of complex numbers of the form \(\exp(2i\pi l/(a_{0}a_{1}\cdots a_{r}))\) where \(l \in Z\) and \(r \in N\) , equipped with the discrete topology. Then \(\chi \mapsto \chi(1)\) is an isomorphism from \(\widehat{\Delta}_{\boldsymbol{a}}\) onto \(\mathbf{Z}(\boldsymbol{a}^{\infty})\) .
\item
  Let B be the subgroup of \(\mathbf{R} \times \Delta_{\boldsymbol{a}}\) consisting of the \((n, n)\) for \(n \in \mathbf{Z}\) ; it is discrete. We set \(\Sigma_{\boldsymbol{a}} = (\mathbf{R} \times \Delta_{\boldsymbol{a}})/\mathrm{B}\) . It is a compact, metrizable, connected group. (Note that the canonical image of \(\mathbf{R}\) into \(\Sigma_{\boldsymbol{a}}\) is dense, and use Lemma 1 of II, p.~244.)
\end{enumerate}

\(d)\) Let \(\Gamma_{\boldsymbol{a}}\) be the additive subgroup of the discrete group \(\mathbf{Q}\) consisting of rational numbers of the form \(l/(a_{0}a_{1}\cdots a_{r})\) where \(l\in\mathbf{Z}\) and \(r\in\mathbf{N}\) . If \(\chi\in\widehat{\Sigma}_{\boldsymbol{a}}\) , then \(\chi\) defines a unitary character of \(\mathbf{R}\times\Delta_{\boldsymbol{a}}\) , hence a unitary character of \(\mathbf{R}\) of the form \(x\mapsto\exp(2i\pi\alpha_{\chi}x)\) where \(\alpha_{\chi}\in\mathbf{R}\) . Then \(\chi\mapsto\alpha_{\chi}\) is an isomorphism from \(\widehat{\Sigma}_{\boldsymbol{a}}\) onto \(\Gamma_{\boldsymbol{a}}\) .

\begin{enumerate}
\def\labelenumi{\alph{enumi})}
\setcounter{enumi}{4}
\item
  Show that \(\widehat{\mathbf{Q}}\) is canonically isomorphic to \(\Sigma_{\pmb{a}}\) , where the sequence \(\pmb{a}\) is defined by \(a_{n} = n + 2\) for all \(n \geqslant 0\) .
\item
  Let \(R_{d}\) be the group R equipped with the discrete topology. Show that \(\widehat{R}_{d}\) is isomorphic to \((\widehat{\Sigma}_{\boldsymbol{a}})^{\mathfrak{c}}\) where \(a_{n}=n+2\) and where c is the cardinality of the continuum. (Consider a basis of R as a Q-vector space.)
\end{enumerate}

¶ 10) A topological group is said to be monothetic if there exists an element of the group whose powers are dense in the group, and solenoidal if there exists a continuous homomorphism from R into the group whose image is dense.

\begin{enumerate}
\def\labelenumi{\alph{enumi})}
\item
  With the notation of Exercise 9, the group \(\Delta_{\pmb{a}}\) is monothetic.
\item
  Every compact monothetic totally disconnected group is topologically isomorphic to a group of the form \(\Delta_{a}\) .
\item
  Suppose G is compact. For G to be monothetic, it is necessary and sufficient that \(\widehat{G}\) is isomorphic to a subgroup of T equipped with the discrete topology.
\item
  Suppose G is compact. The following conditions are equivalent:
\end{enumerate}

\begin{enumerate}
\def\labelenumi{(\roman{enumi})}
\item
  G is solenoidal;
\item
  \(\widehat{\mathbf{G}}\) is isomorphic to a subgroup of \(\mathbf{R}\) endowed with the discrete topology;
\item
  \(\widehat{\mathbf{G}}\) is torsion-free, and \(\operatorname{Card}(\widehat{\mathbf{G}}) \leqslant \mathfrak{c}\), the cardinality of the continuum;
\item
  G is a quotient of \((\Sigma_{\boldsymbol{a}})^{\mathfrak{c}}\) with \(\boldsymbol{a}=(n+2)_{n\geqslant0}\) .
\end{enumerate}

\begin{enumerate}
\def\labelenumi{\alph{enumi})}
\setcounter{enumi}{4}
\tightlist
\item
  If G is a compact solenoidal group, then G is monothetic.
\end{enumerate}

\begin{enumerate}
\def\labelenumi{\arabic{enumi})}
\setcounter{enumi}{10}
\tightlist
\item
  We denote by L(G) the set of continuous representations of G into R. A one-parameter subgroup of G is the image of a continuous homomorphism from R into G.
\end{enumerate}

\begin{enumerate}
\def\labelenumi{\alph{enumi})}
\tightlist
\item
  The following conditions are equivalent:
\end{enumerate}

\begin{enumerate}
\def\labelenumi{(\roman{enumi})}
\item
  G is the union of one-parameter subgroups;
\item
  every homomorphism from Z into G extends to a continuous homomorphism from R into G;
\item
  every continuous homomorphism from \(\widehat{\mathbf{G}}\) into \(\mathbf{R} / \mathbf{Z}\) arises by passing to the quotient from an element of \(\mathrm{L}(\widehat{\mathrm{G}})\) ;
\item
  every unitary character of \(\widehat{\mathbf{G}}\) is of the form \(e^{i\lambda}\) where \(\lambda \in \mathrm{L}(\widehat{\mathbf{G}})\) .
\end{enumerate}

\begin{enumerate}
\def\labelenumi{\alph{enumi})}
\setcounter{enumi}{1}
\item
  If G has a countable base, the conditions of a) are equivalent to: G is isomorphic to \(R^{n} \times T^{I}\) where I is a countable set.
\item
  The union of the one-parameter subgroups of G is a dense subgroup of the neutral component of G. (Use Exerc. 4.)
\end{enumerate}

\(d)\) Let \(t \mapsto \chi_t\) be a continuous map from \([0,1]\) into \(\widehat{\mathbf{G}}\) such that \(\chi_0\) is the trivial character. There exists a map \(t \mapsto \lambda_t\) from \([0,1]\) into \(\mathrm{L}(\mathrm{G})\) with \(\lambda_0 = 0\) such that, for all \(t\), we have \(\chi_t = \exp(2i\pi\lambda_t)\) , and, for every \(x \in \mathrm{G}\) , the map \(t \mapsto \lambda_t(x)\) is continuous.

\begin{enumerate}
\def\labelenumi{\alph{enumi})}
\setcounter{enumi}{4}
\tightlist
\item
  The union of the one-parameter subgroups of G is also the union of the arcs of G that contain e, an arc of G being the image of a continuous map from \([0,1]\) into G. (Use d).
\end{enumerate}

\begin{enumerate}
\def\labelenumi{\arabic{enumi})}
\setcounter{enumi}{11}
\tightlist
\item
  We use the notation L(G) from Exercise 11.
\end{enumerate}

\begin{enumerate}
\def\labelenumi{\alph{enumi})}
\item
  Let H be a closed subgroup of G. Every element of L(H) extends to an element of L(G). (First reduce to the case where \(G = R^{n} \times D\) with D discrete, then in the case where the given element of L(H) is trivial on \(R^{n}\) and then to the case where G is discrete. Then use the fact that R is divisible.)
\item
  Every continuous morphism from H into \(R^{n} \times T^{I}\) , where I is an arbitrary set, extends to a continuous morphism from G into \(R^{n} \times T^{I}\) (Use a) and Th. 4 of II, p.~226.)
\item
  Conversely, let A be a locally compact commutative group. Suppose that, for every locally compact commutative group G, for every closed subgroup H of G, and every continuous morphism \(\varphi\) of H into A, \(\varphi\) extends to a continuous morphism from G into A. Then there exist an integer n and a set I such that A is isomorphic to \(R^{n} \times T^{I}\) (Taking G = R, H = Z, one sees that A is connected, hence of the form \(R^{n} \times \widehat{D}\) with D discrete. Show that D is a projective Z-module, hence free.)
\end{enumerate}

\begin{enumerate}
\def\labelenumi{\arabic{enumi})}
\setcounter{enumi}{12}
\tightlist
\item
  Let C be a compact subset of G and \(\varepsilon > 0\) .
\end{enumerate}

\begin{enumerate}
\def\labelenumi{\alph{enumi})}
\item
  Let \(\alpha\) be the Haar measure of G. There exists a compact subset D of G such that \(\alpha(\mathrm{CD}) \leqslant (1 + \varepsilon)\alpha(\mathrm{D})\) . (The case where \(\mathbf{G} = \mathbf{R}^n\) being immediate, we are reduced to the case where G admits a compact open subgroup H, and we may assume C is saturated with respect to H; we are then reduced to the case where G is discrete; we may next assume G is finitely generated, and finally \(\mathbf{G} = \mathbf{Z}^p\) .)
\item
  There exists \(k \in \mathrm{L}^{1}(\widehat{\mathrm{G}})\) such that \(\mathcal{F}(k) \in \mathcal{K}(\mathrm{G})\) and \(\mathcal{F}(k)(x) = 1\) for \(x \in C\) , with \(\|k\|_{1} \leqslant 1 + \varepsilon\) . (Take k of the form \(\lambda fg\) , where \(\lambda\) is a constant and where \(f, g \in \mathrm{L}^{2}(\widehat{\mathrm{G}})\) have as Fourier transforms the characteristic functions of suitable sets, using a).)
\end{enumerate}

¶ 14) a) Suppose that every neighborhood of e in G contains an element of infinite order. There exists a compact metrizable perfect totally disconnected subset P of G which is a Kronecker set (Exercise 15 of II, p.~265), and a nonzero diffuse measure concentrated on P. (Reduce to the case where G is metrizable using Exercise 3, c). Then imitate the construction of TG, IX, p.~64, Lemma 3, and at the same time use part e) of Exercise 15 of II, p.~265.)

\begin{enumerate}
\def\labelenumi{\alph{enumi})}
\setcounter{enumi}{1}
\item
  If \(\mathrm{G} = (\mathbf{Z}/q\mathbf{Z})^{\mathbf{N}}\) , there exists a compact perfect totally disconnected subset P of G which is of type \(K_{q}\) , and a nonzero diffuse measure concentrated on P.
\item
  Assume G is non-discrete. There exists an element \(\tau\) of \(\mathcal{M}^{1}(G)\) such that \(\|\mathcal{F}(\tau)\|_{\infty}<\varrho(\tau)\) , where \(\varrho(\tau)\) is the spectral radius in \(\mathcal{M}^{1}(G)\) . There exists a non-invertible element \(\tau'\) of \(\mathcal{M}^{1}(G)\) such that \(\mathcal{F}(\tau')\geqslant1\) everywhere on \(\widehat{G}\) . There exists a non-Hermitian character of \(\mathcal{M}^{1}(G)\) . (Use a), b), exerc. 3, c), and part f) of exerc. 15 of II, p.~265.)
\end{enumerate}

\(d)\) Deduce from \(c\) ) that \(\widehat{\mathrm{G}}\) , which is open in \(\mathsf{X}(\mathcal{M}^1 (\mathrm{G}))\) , is not dense in \(\mathsf{X}(\mathcal{M}^1 (\mathrm{G}))\) .

\exercisesection{Invariant Subspaces}

\begin{enumerate}
\def\labelenumi{\arabic{enumi})}
\tightlist
\item
  Let \((a_{n})_{n\in \mathbf{N}}\) be a sequence of complex numbers such that there exists a real number \(A\geqslant 0\) satisfying \(n|a_n|\leqslant A\) for all \(n\in \mathbf{N}\) . For \(z\in \mathbf{C}\) of modulus \(<  1\) , set \(\varphi (z) = \sum_{n\geqslant 0}a_nz^n\) .
\end{enumerate}

One assumes that there exists \(\ell \in \mathbf{C}\) such that

\[
\lim_{\substack{x\in \mathbf{R}_{+}\\ x <   1}}\varphi (x) = \ell .
\]

For \(t\in \mathbf{R}_+^*\) , set

\[
g (t) = \sum_ {0 \leqslant n \leqslant t ^ {- 1}} a _ {n}, \quad \mathrm{F} (t) = \varphi (e ^ {- t}).
\]

We denote by \(dx\) the Lebesgue measure on \(\mathbf{R}\) as well as its restriction to \(\mathbf{R}_{+}^{*}\) .

\begin{enumerate}
\def\labelenumi{\alph{enumi})}
\tightlist
\item
  Let \(s \leqslant t\) be elements of \(\mathbf{R}_{+}^{*}\) . We have
\end{enumerate}

\[
| g (s) - g (t) | \leqslant | a _ {[ t ^ {- 1} + 1 ]} | + A \log \left(\frac {t}{s}\right).
\]

\begin{enumerate}
\def\labelenumi{\alph{enumi})}
\setcounter{enumi}{1}
\item
  The function g is slowly oscillating for \(t \rightarrow 0\) .
\item
  For every \(t \in R_{+}^{*}\) , the function \(x \mapsto g(t/x)xe^{-x}\) is integrable on \(R_{+}^{*}\) with respect to the Haar measure \(x^{-1}dx\) .
\item
  For \(0 < x < 1\) , one has
\end{enumerate}

\[
\varphi (x) = \sum_ {n \geqslant 0} g \left(\frac {1}{n}\right) \left(x ^ {n} - x ^ {n + 1}\right).
\]

\begin{enumerate}
\def\labelenumi{\alph{enumi})}
\setcounter{enumi}{4}
\tightlist
\item
  Let \(f(x) = xe^{-x}\) for \(x \in R_{+}^{*}\) . We have \(f \in \mathrm{L}^{1}(\mathbf{R}_{+}^{*}, dx/x)\) ; its integral is equal to 1, and we have F = g * f.
\end{enumerate}

\(f)\) The function \(g\) belongs to \(\mathrm{L}^{\infty}(\mathbf{R}_{+}^{*})\) . (Show that \(\mathrm{F}-g\) is bounded using \(a\) ) as well as the hypothesis.)

\(g)\) The Fourier transform of the function \(f\) on the group \(\mathbf{R}_{+}^{*}\) does not vanish. (Identify \(\mathcal{F}(f)\) to the function \(y\mapsto\Gamma(1+iy)\) and use FVR, VII.)

\begin{enumerate}
\def\labelenumi{\alph{enumi})}
\setcounter{enumi}{7}
\tightlist
\item
  We have
\end{enumerate}

\[
\sum_ {n = 0} ^ {+ \infty} a _ {n} = \ell .
\]

This statement is Littlewood's "Tauberian" theorem. The hypothesis concerning the sequence \((a_n)\) , that is, \(a_n = \mathrm{O}(1/n)\) , can moreover be weakened: it suffices that there exist a constant \(A \geqslant 0\) such that \(na_n \geqslant -A\) for every integer \(n \geqslant 1\) (“Hardy–Littlewood Tauberian theorem,” see for example J. KOREVAAR, Tauberian theory, Grundlehren math. wiss. 329, Springer (2004), Th. 7.4, p.~16; the reader should take care not to confuse this theorem with that of FVR, I, p.~50, Exercise 18).

\begin{enumerate}
\def\labelenumi{\arabic{enumi})}
\setcounter{enumi}{1}
\item
  Let F be a closed subset of \(\widehat{\mathbf{G}}\) . Let K be a compact subset of \(\widehat{\mathbf{G}}\) such that \(\mathrm{F} \cap \mathrm{K} = \varnothing\) . Show that there exists a continuous function \(g\) integrable on G such that \(\mathcal{F}(g)\) is equal to 0 on F and to 1 on K. (Let \(f \in \mathrm{L}^{1}(\mathrm{G})\) be such that \(\mathcal{F}(f)\) is zero on F and equal to 1 on K; consider \(f * \varphi\) , where \(\varphi\) is the Fourier transform of a function of the form \(h * h'\) with \(h \in \mathcal{K}(\widehat{\mathrm{G}})\) , \(h' \in \mathcal{K}(\widehat{\mathrm{G}})\) , and where \(h * h' = 1\) on K.)
\item
  Let \(g\) be a complex function defined on \(]0, +\infty[\) , integrable with respect to Lebesgue measure. We assume that \(\int_0^{+\infty} g(t)dt = 1\) and that \(\int_0^{+\infty} g(t)t^{ix}dt \neq 0\) for all \(x \in \mathbf{R}\) . Let \(f\) be a complex measurable and bounded function in \(]0, +\infty[\) , such that \(\frac{1}{x}\int_0^{+\infty} g(t/x)f(t)dt\) tends to a finite limit \(\ell\) when \(x\) tends to \(+\infty\) .
\end{enumerate}

\begin{enumerate}
\def\labelenumi{\alph{enumi})}
\tightlist
\item
  For any complex-valued function h defined on \(]0, +\infty[\) , integrable with respect to Lebesgue measure and with integral 1, we have
\end{enumerate}

\[
\lim _ {x \to + \infty} \frac {1}{x} \int_ {0} ^ {+ \infty} h \left(\frac {t}{x}\right) f (t) d t = \ell .
\]

In particular,

\[
\lim _ {x \rightarrow + \infty} \frac {1}{x} \int_ {0} ^ {x} f (t) d t = \ell .
\]

\((\text{Setting } g_1(t) = tg(t), h_1(t) = th(t), \text{ we have } g_1 \in \mathrm{L}^1(\mathbf{R}_+^*) ; \text{ the Fourier transform } \mathcal{F}(g_1) \text{ does not vanish and } \check{g}_1 * f \text{ tends to } \ell, \text{ hence } h_1 * f \text{ tends to } \ell.)\) b) If \(f(t)\) is slowly oscillating on \(\mathbf{R}_+^*\) when \(t\) tends to \(+\infty\) , then \(f\) tends to \(\ell\) when \(t\) tends to \(+\infty\) .

\begin{enumerate}
\def\labelenumi{\arabic{enumi})}
\setcounter{enumi}{3}
\tightlist
\item
  Let G = R. For every real number \(\alpha > 1\) , we define \(f_{\alpha} \in L^{1}(\mathbf{R})\) by
\end{enumerate}

\[
f _ {\alpha} (x) = \left\{ \begin{array}{l l} 2 & \text { if } 0 <   x <   1 \\ 1 & \text { if } 1 \leqslant x <   \alpha \\ 0 & \text { otherwise. } \end{array} \right.
\]

\begin{enumerate}
\def\labelenumi{\alph{enumi})}
\item
  The functions \(f * \varepsilon_{x}\) for \(x \in R\) form a total set in \(\mathrm{L}^{1}(\mathbf{R})\) if and only if \(\alpha\) is irrational.
\item
  The functions \(f * \varepsilon_x\) for \(x \in \mathbf{R}\) form a total set in \(L^2(\mathbf{R})\) .
\end{enumerate}

\begin{enumerate}
\def\labelenumi{\arabic{enumi})}
\setcounter{enumi}{4}
\tightlist
\item
  Let \(f \in \mathrm{L}^{\infty}([0, +\infty([)\text{ and } k \in \mathbf{R}_{+}^{*}. \text{ If}\)
\end{enumerate}

\[
\lim _ {x \to + \infty} \frac {1}{x} \int_ {0} ^ {+ \infty} e ^ {- t / x} f (t) d t = \ell
\]

then

\[
\lim _ {x \rightarrow + \infty} \frac {k}{x ^ {k}} \int_ {0} ^ {x} (x - t) ^ {k - 1} f (t) d t = \ell
\]

(Use the functions \(g(t) = e^{-t}\) , and \(h(t) = k(1 - t)^{k-1}\) for \(0 \leqslant t \leqslant 1\) , \(h(t) = 0\) for \(t > 1\) .)

\begin{enumerate}
\def\labelenumi{\arabic{enumi})}
\setcounter{enumi}{5}
\tightlist
\item
  Let \(f \colon \mathbf{R} \to \mathbf{R}\) and \(g: \mathbf{R} \to \mathbf{C}\) be two functions. We assume that \(f \in \mathrm{L}^{\infty}(\mathbf{R})\) and that \(g \in \mathrm{L}^{1}(\mathbf{R})\) , \(\int_{\mathbf{R}} g(x) dx = 1\) and that \(\mathcal{F}(g)\) does not vanish. We assume that \(f\) is slowly decreasing, that is, for \(y > x\) such that \(y - x\) tends to 0 and \(x \to +\infty\) , one has
\end{enumerate}

\[
\liminf (f (y) - f (x)) \geqslant 0.
\]

If \((g*f)(x)\) tends to \(\ell\) as x tends to \(+\infty\) , then \(f(x)\) tends to \(\ell\) as x tends to \(+\infty\) .

\begin{enumerate}
\def\labelenumi{\arabic{enumi})}
\setcounter{enumi}{6}
\tightlist
\item
  Let \(g: R \rightarrow C\) be a continuous function such that:
\end{enumerate}

\[
\sum_ {n = - \infty} ^ {+ \infty} \sup _ {n \leqslant t \leqslant n + 1} | g (t) | <   + \infty .
\]

\begin{enumerate}
\def\labelenumi{\alph{enumi})}
\item
  We have \(g \in \mathrm{L}^{1}(\mathbf{R})\) .
\item
  Suppose that g has integral 1 and that \(\mathcal{F}(g)\) does not vanish. Let \(\mu\) be a complex measure on R such that \(|\mu|([x,x+1])\) is bounded as x ranges over R. If \((g*\mu)(x)\) tends to \(\ell\) as x tends to \(+\infty\) then, for every continuous function \(h:R\to C\) such that
\end{enumerate}

\[
\sum_ {n = - \infty} ^ {+ \infty} \sup _ {n \leqslant t \leqslant n + 1} | h (t) | <   + \infty ,
\]

of integral 1, the limit of \((h * \mu)(x)\) when \(x\) tends to \(+\infty\) is equal to \(\ell\) . (Show that \(h * \mu\) is slowly oscillating and that \((g * (h * \mu))(x)\) tends to \(\ell\) when \(x\) tends to \(+\infty\) .)

\begin{enumerate}
\def\labelenumi{\arabic{enumi})}
\setcounter{enumi}{7}
\tightlist
\item
  Let \(g\) be a function satisfying the hypotheses of Exerc. 3. Moreover, suppose \(g \geqslant 0\) and \(\liminf_{x \to 0} g(x) > 0\) . Let \(f\) be a real-valued measurable function bounded below on \(]0, +\infty[\) . If
\end{enumerate}

\[
\lim _ {x \to + \infty} \frac {1}{x} \int_ {0} ^ {+ \infty} g \left(\frac {t}{x}\right) f (t) d t = \ell
\]

then

\[
\sigma (x) = \frac {1}{x} \int_ {0} ^ {x} f (t) d t
\]

tends to \(\ell\) when \(x\) tends to \(+\infty\) . (Show that \(\sigma\) is bounded for \(x \geqslant 1\) , then that \(\sigma\) is slowly decreasing, then that

\[
\frac {1}{x} \int_ {0} ^ {+ \infty} g \left(\frac {t}{x}\right) \sigma (t) d t
\]

tends to \(\ell\) when \(x\) tends to \(+\infty\) .)

\begin{enumerate}
\def\labelenumi{\arabic{enumi})}
\setcounter{enumi}{8}
\tightlist
\item
  Let \(\varphi\colon[0,+\infty[\to\mathbf{R}\) be a nonnegative increasing function, such that the function \(t\mapsto e^{-\sigma t}\varphi(t)\) is integrable with respect to Lebesgue measure when \(\sigma>1\) . For every \(s\in\mathbf{C}\) such that \(\mathcal{R}(s)>1\) , set
\end{enumerate}

\[
f (s) = \int_ {0} ^ {+ \infty} e ^ {- s t} \varphi (t) d t
\]

(“Laplace transform of \(\varphi\)”).

Suppose there exists \(A \in C\) with the following property: as \(\sigma\) tends to 1 through values \textgreater{} 1, the function

\[
t \mapsto f (\sigma + i \tau) - \frac {\mathrm{A}}{\sigma + i \tau - 1} \qquad \qquad \tau \in \mathbf {R}
\]

converges uniformly to a function g on every compact subset of R. Then

\[
\lim _ {t \to + \infty} \varphi (t) e ^ {- t} = A
\]

(“Ikehara tauberian theorem”).

(We set \(a(t) = e^{-t}\varphi(t)\) for \(t > 0\) , \(a(t) = 0\) for \(t \leqslant 0\) , \(A(t) = A\) for \(t > 0\) , \(A(t) = 0\) for \(t \leqslant 0\) . Let \(\lambda > 0\). We set:

\[
k _ {\lambda} (t) = \lambda \sqrt {\frac {2}{\pi}} \Big (\frac {\sin 2 \lambda t}{2 \lambda t} \Big) ^ {2}
\]

\[
\mathrm{K} _ {\lambda} (t) = \left\{ \begin{array}{l l} 1 - \left| \frac {t}{2 \lambda} \right| & \text { if } | t | \leqslant 2 \lambda \\ 0 & \text { if } | t | > 2 \lambda . \end{array} \right.
\]

Using exercise 1 of II, p.~262, the Lebesgue–Fubini theorem, and the Plancherel theorem, show that:

\[
\lim _ {\varepsilon \to 0} \frac {1}{\sqrt {2 \pi}} \int_ {- \infty} ^ {+ \infty} k _ {\lambda} (x - t) (a (t) - \mathrm{A} (t)) e ^ {- \varepsilon t} d t = \frac {1}{\sqrt {2 \pi}} \int_ {- 2 \lambda} ^ {2 \lambda} \mathrm{K} _ {\lambda} (y) e ^ {- i x y} g (y) d y
\]

and deduce that

\[
\lim _ {x \to + \infty} \frac {1}{\sqrt {2 \pi}} \int_ {- \infty} ^ {+ \infty} k _ {\lambda} (x - t) a (t) d t = \mathrm{A}.
\]

Prove then that \(a\) is bounded and slowly decreasing. One can finally apply exerc. 6.)

\begin{enumerate}
\def\labelenumi{\arabic{enumi})}
\setcounter{enumi}{9}
\tightlist
\item
  We define the von Mangoldt function on \(\mathbf{N}^{*}\) by \(\Lambda(n)=0\) if n is not of the form \(p^{k}\) where p is prime and \(k\geqslant1\) , and \(\Lambda(p^{k})=\log(p)\) .
\end{enumerate}

We denote by \(\zeta\) the holomorphic function on \(\mathbf{C} - \{1\}\) whose restriction to the set of \(s \in \mathbf{C}\) such that \(\mathcal{R}(s) > 1\) coincides with the Riemann zeta function (Exercise 16 of II, p.~266).

For \(x > 0\), we set

\[
\psi (x) = \sum_ {1 \leqslant n \leqslant x} \Lambda (n).
\]

\begin{enumerate}
\def\labelenumi{\alph{enumi})}
\tightlist
\item
  Show that for all \(s \in \mathbf{C}\) such that \(\mathcal{R}(s) > 1\) , one has
\end{enumerate}

\[
\sum_ {n \geqslant 1} \Lambda (n) n ^ {- s} = - \frac {\zeta^ {\prime} (s)}{\zeta (s)}.
\]

\begin{enumerate}
\def\labelenumi{\alph{enumi})}
\setcounter{enumi}{1}
\tightlist
\item
  We have
\end{enumerate}

\[
\lim _ {x \to + \infty} \frac {\psi (x)}{x} = 1.
\]

(Use Exercise 16 of II, p.~266 and Exercise 9).

\begin{enumerate}
\def\labelenumi{\alph{enumi})}
\setcounter{enumi}{2}
\tightlist
\item
  Let \(\pi(x)\) be the number of prime numbers \(\leqslant x\) . Show that
\end{enumerate}

\[
\pi (x) \sim \frac {x}{\log (x)}
\]

when \(x \to +\infty\) (“prime number theorem,” due to Hadamard and de la Vallée Poussin).

\begin{enumerate}
\def\labelenumi{\arabic{enumi})}
\setcounter{enumi}{10}
\tightlist
\item
  One defines the Möbius function \(\mu\) on \(\mathbf{N}^*\) by \(\mu(n) = (-1)^k\) if \(n\) is the product of \(k\) prime factors pairwise distinct and \(\mu(n) = 0\) otherwise.
\end{enumerate}

For \(x > 0\), we set

\[
\mathrm{M} (x) = \sum_ {n \leqslant x} \mu (n), \quad f (x) = \frac {\mathrm{M} (x)}{x}.
\]

\begin{enumerate}
\def\labelenumi{\alph{enumi})}
\tightlist
\item
  Show that for \(\mathcal{R}(s) > 1\) , one has
\end{enumerate}

\[
\sum_ {n \geqslant 1} \mu (n) n ^ {- s} = \frac {1}{\zeta (s)}.
\]

\begin{enumerate}
\def\labelenumi{\alph{enumi})}
\setcounter{enumi}{1}
\item
  Show that f is bounded, and that \(f(y) - f(x) = \mathrm{O}((y - x)/x)\) for \(y \geqslant x \geqslant 1\) , so that f is slowly oscillating on the group \(R_{+}^{*}\) as x tends to \(+\infty\) .
\item
  Let \(a, b \in R_{+}^{*}\) . For \(x \textgreater{} 0\), let {[}x{]} be the integer part of x. We set \(g_{0}(t) = [t^{-1}]\) for \(t \in R_{+}^{*}\) and
\end{enumerate}

\[
g (t) = 2 g _ {0} (t) - a g _ {0} (a t) - b g _ {0} (b t).
\]

The function g is integrable on \(R_{+}^{*}\) with respect to Lebesgue measure. If \(s \in C\) and \(\mathcal{R}(s) > 0\) , one has

\[
\int_ {0} ^ {+ \infty} g (t) t ^ {s} d t = (2 - a ^ {- s} - b ^ {- s}) \frac {\zeta (1 + s)}{1 + s}.
\]

\(d)\) We have \(\int_{0}^{+\infty} g(t)t^{ix} dt \neq 0\) for all \(x \in \mathbf{R}\) . (Use the part \(j\) ) of Exercise 16 of II, p.~266.)

\begin{enumerate}
\def\labelenumi{\alph{enumi})}
\setcounter{enumi}{4}
\tightlist
\item
  Show that \(\int_0^{+\infty} g(t / x)f(t)dt = o(x)\) when \(x \to +\infty\) . Deduce that \(\mathrm{M}(x) = o(x)\) when \(x \to +\infty\) .
\end{enumerate}

(One can show in an elementary way that this result is equivalent to the prime number theorem, cf.~Exercise 10; see for example H. IWANIEC and E. KOWALSKI, Analytic number theory, A.M.S Coll. Publ. 53 (2004), pp.~31--32.)

\begin{enumerate}
\def\labelenumi{\arabic{enumi})}
\setcounter{enumi}{11}
\item
  Let J be the set of \(f \in \mathcal{S}(\mathbf{R}^3)\) such that \(\mathcal{F}(f)\) vanishes on the sphere \(\mathbf{S}_2\) . Let I be the set of \(f \in \mathrm{J}\) such that \((\partial/\partial y_1)(\mathcal{F}(f))\) vanishes on \(\mathbf{S}_2\) . Let \(\overline{\mathrm{I}}\) and \(\overline{\mathrm{J}}\) be the closures of I and J in \(\mathrm{L}^1(\mathbf{R}^3)\) . Then we have \(\mathrm{V}(\overline{\mathrm{I}}) = \mathrm{V}(\overline{\mathrm{J}}) = \mathbf{S}_2\) but \(\overline{\mathrm{I}} \neq \overline{\mathrm{J}}\) . (Let \(\mu\) be the positive measure of mass 1 on \(\mathbf{S}_2\) invariant under the orthogonal group of \(\mathbf{R}^3\) ; show, using Exercise 20 of II, p.~269, c), that \(f(t) = t_1\mathcal{F}(\mu)(t)\) is an element of \(\mathrm{L}^\infty(\mathbf{R}^3)\) orthogonal to I but not to J.)
\item
  One calls a C-set in \(\widehat{\mathbf{G}}\) a closed subset E of \(\widehat{\mathbf{G}}\) having the following property: if \(f\in \mathrm{L}^{1}(\mathrm{G})\) , if \(\mathcal{F}(f)\) vanishes on E, and if \(\varepsilon >0\) , there exists a function \(g\in \mathrm{L}^{1}(\mathrm{G})\) such that \(\| f - f*g\| \leqslant \varepsilon\) and such that \(\operatorname {Supp}(\mathcal{F}(g))\) be compact and disjoint from E.
\end{enumerate}

\begin{enumerate}
\def\labelenumi{\alph{enumi})}
\item
  If E is a C-set in \(\widehat{G}\) , there exists only one closed ideal I of \(\mathrm{L}^1 (\mathrm{G})\) such that \(\mathrm{V(I)} = \mathrm{E}\) .
\item
  Every singleton subset of \(\widehat{\mathrm{G}}\) is a C-set.
\item
  The union of two C-sets is a C-set.
\item
  Every set obtained by translation from a C-set is a C-set.
\item
  Let H be a closed subgroup of \(\widehat{\mathbf{G}}\) , let E be a closed subset of H, and let E' be the boundary of E relative to H. If E' is a C-set relative to \(\widehat{\mathbf{G}}\) , then E is a C-set with respect to \(\widehat{\mathbf{G}}\) . (Use exerc. 21 of II, p.~270.)
\item
  In \(R^{n}\) , every intersection E of a finite number of closed half-spaces is a C-set. (Reason by induction on the dimension of the affine subspace generated by E, using b), c), d), e).
\end{enumerate}

\begin{enumerate}
\def\labelenumi{\arabic{enumi})}
\setcounter{enumi}{13}
\tightlist
\item
  Let E be a closed subset of \(R^{n}\) . Suppose that there exists an interior point p of E such that every line passing through p meets the boundary of E in at most two points. Then there exists only one closed ideal I of \(L^{1}(\mathbf{R}^{n})\) such that \(\mathcal{F}(f)\) vanishes on E. (Consider f as a limit in \(L^{1}(\mathbf{R}^{n})\) of functions deduced from f by homotheties with center p).
\end{enumerate}

¶ 15) We denote by μ a Haar measure of G and by B the filter of neighborhoods of e in \(\widehat{G}\) .

\begin{enumerate}
\def\labelenumi{\alph{enumi})}
\tightlist
\item
  Let U be an integrable open neighborhood of e in G. There exists a function \(a \in \mathrm{L}^{1}(\mathrm{G})\) , \(a \geqslant 0\) , of integral 1, vanishing outside U, such that \(\mathcal{F}(a) \in \mathrm{L}^{1}(\widehat{\mathrm{G}})\) , and
\end{enumerate}

\[
\int_ {\mathrm{G}} a (x) ^ {2} d x \leqslant \frac {2}{\mu (\mathrm{U})}.
\]

(Let V be a compact neighborhood of e such that \(V \subset U\) , \(\mu(U) < \sqrt{2}\mu(V)\) ; let W be a compact symmetric neighborhood of e such that \(VW^{2} \subset U\) ; take \(a = \lambda f * g\) , where \(\lambda\) is a constant \textgreater{} 0, and where f, g are the characteristic functions of the sets VW and W.)

\begin{enumerate}
\def\labelenumi{\alph{enumi})}
\setcounter{enumi}{1}
\tightlist
\item
  If \(f \in \mathrm{L}^{\infty}(\mathrm{G})\) , we denote by \(\mathrm{A}(f)\) the set of elements of \(\widehat{G}\) which belong to the weakly closed vector subspace of \(\mathrm{L}^{\infty}(\mathrm{G})\) invariant under translation generated by f. If \(f, g \in \mathrm{L}^{\infty}(\mathrm{G})\) , one has
\end{enumerate}

\[
\mathrm{A} (f g) \subset \overline {{\mathrm{A} (f) \mathrm{A} (g)}}.
\]

(Let U be a neighborhood of \(e\) . Show that \(fg\) is the weak limit of linear combinations of elements of \(\widehat{G}\) belonging to \(\mathrm{A}(f)\mathrm{UA}(g)\mathrm{U}\) .)

\begin{enumerate}
\def\labelenumi{\alph{enumi})}
\setcounter{enumi}{2}
\tightlist
\item
  Let \(g \in \mathrm{L}^{\infty}(\mathrm{G})\) , K a compact subset of \(\widehat{G}\) , and f a function in \(\mathrm{L}^{1}(\mathrm{G})\) such that \(\mathcal{F}(f)\) vanishes on \(\mathrm{A} = \mathrm{A}(g)\) and outside K. For every \(V \in B\) , let \(\omega(V)\) be the supremum of \(|\mathcal{F}(f)|\) on AV. If
\end{enumerate}

\[
\operatorname * {l i m i n f} _ {\mathfrak {B}} \frac {\omega (\mathrm{V}) ^ {2} \mu ((\mathrm{AV} - \mathrm{A}) \cap \mathrm{K})}{\mu (\mathrm{V})} = 0,
\]

then \(\langle f,g\rangle=0\) . (Let \(\varepsilon>0\) . There exists a compact subset H of G such that \(\int_{G-H}|f|dx\leqslant\varepsilon\) . Then there exists an open neighborhood U of e in \(\widehat{G}\) such that \(|\langle x,\widehat{x}\rangle-1|\leqslant\varepsilon\) for \(x\in H\) , \(\widehat{x}\in U\) . Applying a) to \(\widehat{G}\) and U, one obtains a function \(a\in L^{1}(\widehat{G})\) . Let \(b=\mathcal{F}(a)\in L^{1}(G)\) . Using b), show that \(\mathcal{F}(bg)\) vanishes outside AU. On the other hand,

\[
\left| \int_ {\mathrm{G}} f g d x - \int_ {\mathrm{G}} f g b d x \right| \leqslant \varepsilon (\| f \| _ {1} + 2) \| g \| _ {\infty}
\]

and, since \(f, gb \in \mathrm{L}^{2}(\mathrm{G})\) ,

\[
\int_ {\mathrm{G}} f g b d x = \int_ {\widehat {\mathrm{G}}} \mathcal {F} (f) \mathcal {F} (g b) d \widehat {x} = \int_ {(\mathrm{AU-A}) \cap \mathrm{K}} \mathcal {F} (f) \mathcal {F} (g b) d \widehat {x}.
\]

Bound the absolute value of this last expression by:

\[
2 \| g \| _ {\infty} \frac {\omega (\mathrm{U}) \sqrt {\mu ((\mathrm{AU} - \mathrm{A}) \cap \mathrm{K})}}{\sqrt {\mu (\mathrm{U})}}
\]

to conclude.)

\(d)\) Let I be a closed ideal of \(\mathrm{L}^1 (\mathrm{G})\) , \(\mathrm{A} = \mathrm{V}(\mathrm{I})\) , \(f\) a function of \(\mathrm{L}^1 (\mathrm{G})\) such that \(\mathcal{F}(f)\) vanishes on A, S the set of points where \(\mathcal{F}(f)\) does not vanish, \(\mathrm{S}'\) the boundary of S, and B the largest perfect set contained in \(\mathrm{S} \cap \mathrm{A}\) . For \(\mathrm{V} \in \mathfrak{B}\) , let \(\omega (\mathrm{V})\) be the supremum of \(|\mathcal{F}(f)|\) on AV. We assume that every point of B has a neighborhood K in \(\widehat{\mathrm{G}}\) such that

\[
\operatorname * {l i m i n f} _ {\mathfrak {B}} \frac {\omega (\mathrm{V}) ^ {2} \mu ((\mathrm{AV-A}) \cap \mathrm{K})}{\mu (\mathrm{V})} = 0.
\]

Then \(f \in I\) . (Use the fact that \(L^1(G)\) satisfies Ditkin's condition, and argue as in Prop. 5 of I, p.~94; let \(G' = \widehat{G} \cup \{\infty\}\) be the Alexandroff compactification of \(\widehat{G}\) and let N be the set of \(\chi \in G'\) such that \(\mathcal{F}(f)\) does not belong to \(\mathcal{F}(I)\) in a neighborhood of \(\chi\) . First show that \(N \subset B \cup \{\infty\}\) . Then, using \(c\) ), show that \(N \subset \{\infty\}\) . Finally, show that \(N = \varnothing\) .)

\begin{enumerate}
\def\labelenumi{\alph{enumi})}
\setcounter{enumi}{4}
\tightlist
\item
  Let I be a closed ideal of \(\mathrm{L}^{1}(\mathbf{R}^{n})\) , A = V(I), and f a function in \(\mathrm{L}^{1}(\mathbf{R}^{n})\) such that \(\mathcal{F}(f)\) vanishes on A. For h \textgreater{} 0, let \(A_{h}\) be the set of points of \(R^{n}\) outside A and located at a distance from A less than h, and let \(\omega(h) = \sup_{x \in \mathrm{A}_{h}} |\mathcal{F}(f)(x)|\) . If
\end{enumerate}

\[
\liminf _ {h \to 0} \frac {\omega (h) ^ {2} \mu (\mathrm{A} _ {h})}{h ^ {n}} = 0
\]

then \(f\in \mathrm{I}\) . (Use \(d\) ).)

\(f)\) Let \(f \in \mathrm{L}^{1}(\mathbf{R})\) , \(\alpha \in]0,1[\) , and suppose \(\int_{\mathbf{R}} |y|^{\alpha}| f(y)| dy < +\infty\) . Then there exists a constant \(k\) such that \(|\mathcal{F}(f)(x+h)-\mathcal{F}(f)(x)| \leqslant kh^{\alpha}\) for all \(x \in \mathbf{R}\) and every \(h \in \mathbf{R}\) .

\(g)\) Let I be a closed ideal of \(\mathrm{L}^1 (\mathbf{R})\) , \(f\) a function of \(\mathrm{L}^1 (\mathbf{R})\) such that \(\mathcal{F}(f)\) vanishes on \(\mathrm{V(I)}\) . If

\[
\int_ {\mathbf {R}} | y ^ {1 / 2} f (y) | d y <   + \infty ,
\]

then we have \(f \in I\) . (Use parts e) and f).

\begin{enumerate}
\def\labelenumi{\arabic{enumi})}
\setcounter{enumi}{15}
\tightlist
\item
  Let \(f \in \mathrm{L}^{\infty}(\mathrm{G})\) . We say that \(f\) satisfies the spectral synthesis condition if it belongs to \(\mathrm{Y}(\mathrm{A}(\mathrm{W}_f))\) , where \(\mathrm{W}_f\) is the intersection of all translation-invariant weakly closed subspaces of \(\mathrm{L}^{\infty}(\mathrm{G})\) which contain \(f\) .
\end{enumerate}

\begin{enumerate}
\def\labelenumi{\alph{enumi})}
\item
  Let \(\mu \in \mathcal{M}^1 (\widehat{\mathrm{G}})\) and \(f = \overline{\mathcal{F}}_{\widehat{\mathrm{G}}}(\mu)\in \mathrm{L}^{\infty}(\mathrm{G})\) . Then \(\mathrm{A}(\mathrm{W}_f)\) is the support of \(\mu\) . Moreover, \(f\) satisfies the spectral synthesis condition.
\item
  If G is discrete, every function \(f \in \mathrm{L}^{2}(\mathrm{G})\) satisfies the spectral synthesis condition.
\end{enumerate}

\begin{enumerate}
\def\labelenumi{\arabic{enumi})}
\setcounter{enumi}{16}
\tightlist
\item
  For every \(p \geqslant 1\) , we identify \(\mathrm{L}^p(\mathbf{R}_+)\) with a closed subspace of \(\mathrm{L}^p(\mathbf{R})\) . Let \(f \in \mathrm{L}^1(\mathbf{R}_+)\) be the function defined by \(f(x) = 2e^{-x}\) for \(x \in \mathbf{R}_+\) , and let \(g \in \mathrm{L}^1(\mathbf{R})\) be the function such that \(g(x) = f(-x)\) .
\end{enumerate}

\begin{enumerate}
\def\labelenumi{\alph{enumi})}
\item
  Show that \(f + g = f * g\) .
\item
  Let A (resp. B) be the subalgebra of \(\mathrm{L}^{1}(\mathbf{R})\) generated by f (resp. by f and g). Show that A is dense in \(\mathrm{L}^{1}(\mathbf{R}_{+})\) . (Let \(h \in \mathrm{L}^{\infty}(\mathbf{R}_{+})\) such that \(\langle h, f^{*n} \rangle = 0\) for every integer \(n \geqslant 1\) , where \(f^{*n}\) is the n-fold iterated convolution of f; show that the Laplace transform
\end{enumerate}

\[
\mathrm{F} (z) = \int_ {0} ^ {+ \infty} e ^ {- x z} h (x) d x,
\]

defined for all \(z \in C\) such that \(\mathcal{R}(z) > 0\) is zero.) Deduce that B is dense in \(\mathrm{L}^{1}(\mathbf{R})\) .

\begin{enumerate}
\def\labelenumi{\alph{enumi})}
\setcounter{enumi}{2}
\item
  Let C be a proper closed subalgebra of \(\mathrm{L}^{1}(\mathbf{R})\) containing A. We have \(1 \in \mathrm{Sp}_{\mathrm{C}}(f)\) . (Show that otherwise, part (a) and the holomorphic functional calculus in C would imply that \(g \in C\) .)
\item
  There exists a character \(\chi\) of C such that \(\chi(f)=1\) ; we have
\end{enumerate}

\[
\chi (\varphi) = \int_ {0} ^ {+ \infty} e ^ {- x} \varphi (x) d x
\]

for every function \(\varphi\in\mathrm{L}^{1}(\mathbf{R}_{+})\) .

\begin{enumerate}
\def\labelenumi{\alph{enumi})}
\setcounter{enumi}{4}
\item
  There exists a measure \(\mu \in \mathcal{M}^1 (\mathbf{R})\) of norm 1 such that \(\chi (\varphi) = \int \widehat{\varphi} d\mu\) for every function \(\varphi \in \mathrm{C}\) .
\item
  The Fourier transform of \(\mu\) is the function \(x \mapsto e^{-|x|}\) .
\item
  Let \(\varphi \in \mathrm{C}\) . For every \(y > 0\) , one has
\end{enumerate}

\[
\int_ {\mathbf {R}} \varphi (x) (e ^ {- | x + y |} - e ^ {- | x | - y}) d x = 0.
\]

(Using the relation \(\chi (\varphi *\psi) = \chi (\varphi)\chi (\psi)\) for \(\psi \in \mathrm{L}^1 (\mathbf{R}_+)\) .)

\begin{enumerate}
\def\labelenumi{\alph{enumi})}
\setcounter{enumi}{7}
\tightlist
\item
  We have \(C = L^{1}(\mathbf{R}_{+})\) . (Thus, the subalgebra \(L^{1}(\mathbf{R}_{+})\) is a maximal closed subalgebra of \(L^{1}(\mathbf{R})\) .)
\end{enumerate}

\begin{enumerate}
\def\labelenumi{\arabic{enumi})}
\setcounter{enumi}{17}
\tightlist
\item
  Let G be a discrete ordered commutative topological group. We write the group law of G additively, and denote its identity by 0. We denote by \(G^{+}\) the set of elements \(x \geqslant 0\) of G.
\end{enumerate}

\begin{enumerate}
\def\labelenumi{\alph{enumi})}
\item
  The space \(L_{+}^{1}\) of functions \(f \in \mathrm{L}^{1}(\mathrm{G})\) vanishing outside of \(G^{+}\) is a closed subalgebra of \(\mathrm{L}^{1}(\mathrm{G})\) .
\item
  We assume that G is not Archimedean (TG, V, p.~16, exercise 1 of §3). The algebra L \(_{+}^{1}\) is not maximal in L \(^{1}\) (G).
\end{enumerate}

Assume henceforth that G is Archimedean. Let A be a closed subalgebra of \(\mathrm{L}^{1}(\mathrm{G})\) containing \(L_{+}^{1}\) . For every \(x \in G\) , let \(\varphi_{x} \in \mathrm{L}^{1}(\mathrm{G})\) be the characteristic function of \(\{x\}\) . It is invertible in \(\mathrm{L}^{1}(\mathrm{G})\) and its inverse is \(\varphi_{-x}\) .

\begin{enumerate}
\def\labelenumi{\alph{enumi})}
\setcounter{enumi}{2}
\tightlist
\item
  There exists \(y \in G_{+}\) such that \(\varphi_{y} \in A\) and \(\varphi_{-y} = \varphi_{y}^{-1} \notin A\) and there exists a character \(\chi \neq 0\) of A such that \(\chi(y) = 0\) .
\end{enumerate}

\(d)\) For every \(x > 0\) in G, we have \(\chi(x) = 0\) . (There exists an integer \(k \geqslant 1\) such that \(-y + kx > 0\) .)

\begin{enumerate}
\def\labelenumi{\alph{enumi})}
\setcounter{enumi}{4}
\item
  There exists a measure \(\mu \in \mathcal{M}^1 (\mathrm{G})\) of norm 1 such that \(\chi\) is the restriction of \(\mu\) to the algebra A viewed as a subspace of \(\mathcal{C}(\mathrm{G})\) . (For \(f\in \mathrm{A}\) , one has \(|\chi (f)|\leqslant \varrho (f)\) and moreover \(\varrho (f) = \| f\|_{\infty}\) .)
\item
  Let \(\widehat{\mu}\colon\mathrm{G}\to\mathbf{C}\) be the inverse Fourier transform of \(\mu\) . We have \(\widehat{\mu}(x)=0\) if \(x>0\) .
\item
  The measure \(\mu\) is the counting measure on G.
\item
  The subalgebra \(L_{+}^{1}\) is maximal in \(L^{1}(G)\) .
\end{enumerate}

\begin{enumerate}
\def\labelenumi{\arabic{enumi})}
\setcounter{enumi}{18}
\tightlist
\item
  For \(f \in \mathrm{L}^1(\mathbf{R})\) , we denote by \(\mathrm{A}_f\) the closed subalgebra of \(\mathrm{L}^1(\mathbf{R})\) generated by \(f\) .
\end{enumerate}

\begin{enumerate}
\def\labelenumi{\alph{enumi})}
\tightlist
\item
  There exists \(f \in \mathrm{L}^{1}(\mathbf{R})\) such that \(A_{f}\) is a proper maximal subalgebra of \(\mathrm{L}^{1}(\mathbf{R})\) .
\end{enumerate}

In what follows, we assume that \(f \in \mathrm{L}^{1}(\mathbf{R})\) is a function such that \(A_{f}\) is a proper maximal subalgebra of \(\mathrm{L}^{1}(\mathbf{R})\) .

\begin{enumerate}
\def\labelenumi{\alph{enumi})}
\setcounter{enumi}{1}
\item
  We denote \(S = \{0\} \cup \widehat{f}(\mathbf{R}) \subset \mathbf{C}\) . Show that the function \(z \mapsto \overline{z}\) is not a uniform limit on S of polynomial functions. (First show that C-S is not connected.)
\item
  The subalgebra of \(\mathrm{L}^{1}(\mathbf{R})\) generated by f and \(\widetilde{f}\) is dense in \(\mathrm{L}^{1}(\mathbf{R})\) .
\item
  The Fourier transform of f is injective and does not vanish on R. (Show that S cannot have two bounded connected components.)
\end{enumerate}

\backmatter
\specialchapter{FORMULA SHEET FOR FOURIER THEORY}

\specialsection{Function spaces}

Let G be a locally compact commutative group. We recall the principal function spaces on G and their relation to the Fourier transformation.

The left column defines a function space E on G, and indicates whether it is a subspace of \(L^{1}(G)\) , of \(L^{2}(G)\) , or of both. The right column indicates analogous inclusions for the image of E under the transformation and cotransformation of Fourier (the latter being either that defined on \(L^{1}(G)\) , or that on \(L^{2}(G)\) , as the case may be). The notation F denotes either the Fourier transformation or the Fourier cotransformation.

\(\mathcal{M}^1 (\mathrm{G})\)

\[
\mathcal {F} (\mathcal {M} ^ {1} (\mathrm{G})) \subset \mathcal {C} _ {b} (\widehat {\mathrm{G}})
\]

(space of bounded measures)

(prop. 3, p.~207)

\(\mathrm{L}^{1}(\mathrm{G})\)

\[
\mathcal {F} (\mathrm{L} ^ {1} (\mathrm{G})) \subset \mathcal {C} _ {0} (\widehat {\mathrm{G}})
\]

(integrable functions)

\(\mathrm{L}^2 (\mathrm{G})\)

\[
\mathscr {F} (\mathrm{L} ^ {2} (\mathrm{G})) = \mathrm{L} ^ {2} (\widehat {\mathrm{G}})
\]

(square-integrable functions) (th. 1, p.~215)

\(\mathrm{B(G)}\subset \mathrm{L}^{1}(\mathrm{G})\) \(\mathcal{F}(\mathrm{B(G)}) = \mathrm{B}(\widehat{\mathrm{G}})\)

space of \(f \in \mathrm{L}^1(\mathrm{G})\) such that (th. 3, p.~222)

\[
\mathcal {F} _ {\mathrm{G}} (f) \in \mathrm{L} ^ {1} (\widehat {\mathrm{G}})
\]

\(\mathrm{A(G)}\subset \mathrm{L}^1 (\mathrm{G})\cap \mathrm{L}^2 (\mathrm{G})\cap \mathrm{B(G)}\) \(\mathcal{F}(\mathrm{A(G)})\subset \mathcal{C}_0(\widehat{\mathrm{G}})\cap \mathrm{L}^1 (\widehat{\mathrm{G}})\cap \mathrm{L}^2 (\widehat{\mathrm{G}})\) (space spanned by \(f*g\) where \(f\) , \(g\in \mathrm{L}^{1}(\mathrm{G})\cap \mathrm{L}^{2}(\mathrm{G}))\) (cor. of prop. 8, p.~211, lemma 5, p.~215, prop. 11, p.~217)

\specialsection{\texorpdfstring{Formulary in \(\mathbf{R}^n\)}{Formulary in R\^{}n}}

We recall here the most important formulas concerning the Fourier transform in \(R^{n}\) , equipped with the Euclidean norm. We identify \(R^{n}\) with its dual via the map \((x,y)\mapsto\exp(2i\pi x\cdot y)\) , where \(x\cdot y=\sum_{j}x_{j}y_{j}\) for all \(x=(x_{j})\) and \(y=(y_{j})\) in \(R^{n}\) .

\[
\begin{array}{l l} f \in \mathrm{L} ^ {1} (\mathbf {R} ^ {n}) & \widehat {f} (y) = \int_ {\mathbf {R} ^ {n}} f (x) \exp (- 2 i \pi x \cdot y) d x \\ & (y \in \mathbf {R} ^ {n}; \text {Def. 3, p. 206}) \\ f \in \mathrm{L} ^ {1} (\mathbf {R} ^ {n}) \text {or L} ^ {2} (\mathbf {R} ^ {n}) & \widehat {f} _ {h} (y) = \exp (2 i \pi h \cdot y) \widehat {f} (y) \\ y \in \mathbf {R} ^ {n} & f _ {h} (x) = f (x + h) \\ & (\text {Eq. (11), p. 208}) \\ f \in \mathrm{L} ^ {1} (\mathbf {R} ^ {n}) \text {or L} ^ {2} (\mathbf {R} ^ {n}) & \widehat {g} _ {h} (y) = \widehat {f} (y - h) \\ h \in \mathbf {R} ^ {n} & g _ {h} (x) = f (x) \exp (2 i \pi h \cdot x) \\ & (\text {Eq. (12), p. 208}) \\ f \in \mathrm{L} ^ {1} (\mathbf {R} ^ {n}) \text {or L} ^ {2} (\mathbf {R} ^ {n}) & \widehat {f} \circ \sigma = \frac {1}{| \det (\sigma) |} \widehat {f} \circ {} ^ {t} \sigma^ {- 1} \\ \sigma \in \mathrm{GL} (n, \mathbf {R}) & (\text {Eq. (16), p. 208}) \\ f (x) = \exp (- Q (x)), & \widehat {f} (y) = \frac {1}{\det (\sigma)} \exp (- Q ^ {*} (y)), \\ Q (x) = \| \sigma (x) \| ^ {2} & Q ^ {*} (y) = \| ^ {t} \sigma^ {- 1} (y) \| ^ {2} \\ Q \text {positive definite quadratic form} & (\text {example, p. 239}) \\ f \in \mathrm{L} ^ {2} (\mathbf {R} ^ {n}) & \int_ {\mathbf {R} ^ {n}} | f (x) | ^ {2} d x = \int_ {\mathbf {R} ^ {n}} | \widehat {f} (y) | ^ {2} d y \\ & (\text {Th. 1, p. 215}) \\ f \in \mathrm{L} ^ {2} (\mathbf {R} ^ {n}) \text {such that} & f (x) = \int_ {\mathbf {R} ^ {n}} \widehat {f} (y) \exp (2 i \pi x \cdot y) d x \\ \widehat {f} \in \mathrm{L} ^ {1} (\mathbf {R} ^ {n}) & (\text {almost everywhere ; Prop. 12, p. 219}) \\ f \in A (\mathbf {R} ^ {n}) & f (x) = \int_ {\mathbf {R} ^ {n}} \widehat {f} (y) \exp (2 i \pi x \cdot y) d x \\ & (x \in \mathbf {R} ^ {n}; \text {Prop. 11 p. 217}) \end{array}
\]

\specialsection{\texorpdfstring{Formulary in \((\mathbf{R}/a\mathbf{Z})^{n}\)}{Formulary in (R/aZ)\^{}n}}

For the convenience of the reader, we state the main formulas of Fourier theory for the group \((\mathbf{R}/a\mathbf{Z})^{n}\) , where a \textgreater{} 0, equipped with the Haar measure of total mass 1, denoted \(a^{-n} dx\) . When n = 1, this corresponds to periodic functions of period a. The dual group of \((\mathbf{R}/a\mathbf{Z})^{n}\) is identified with the group \(Z^{n}\) via the map \((\mathbf{R}/a\mathbf{Z})^{n} \times \mathbf{Z}^{n} \to \mathbf{U}\) induced by passing to the quotient from the map defined by

\[
(x, y) \mapsto \exp \left(\frac {2 i \pi}{a} x \cdot y\right)
\]

for all \(x \in R^{n}\) and every \(y \in Z^{n}\) . The dual measure on \(Z^{n}\) is then the counting measure.

\[
\begin{array}{l l} f \in \mathrm{L} ^ {1} ((\mathbf {R} / a \mathbf {Z}) ^ {n}) & \widehat {f} (y) = \frac {1}{a ^ {n}} \int_ {(\mathbf {R} / a \mathbf {Z}) ^ {n}} f (x) \exp \Bigl (- \frac {2 i \pi}{a} y \cdot x \Bigr) d x \\ & (y \in \mathbf {Z} ^ {n}; \text {Def. 3, p. 206}) \\ f \in \mathrm{L} ^ {1} ((\mathbf {R} / a \mathbf {Z}) ^ {n}) \text {or} & \widehat {f} _ {h} (y) = \exp (\frac {2 i \pi}{a} h \cdot y) \widehat {f} (y) \\ \mathrm{L} ^ {2} ((\mathbf {R} / a \mathbf {Z}) ^ {n}) & f _ {h} (x) = f (x + h) \\ h \in (\mathbf {R} / a \mathbf {Z}) ^ {n} & (\text {Eq. (11), p. 208}) \\ f \in \mathrm{L} ^ {1} ((\mathbf {R} / a \mathbf {Z}) ^ {n}) \text {or} & \widehat {g} _ {h} (y) = \widehat {f} (y - h) \\ \mathrm{L} ^ {2} ((\mathbf {R} / a \mathbf {Z}) ^ {n}) & g _ {h} (x) = f (x) \exp (\frac {2 i \pi}{a} h \cdot x) \\ h \in (\mathbf {R} / a \mathbf {Z}) ^ {n} & (\text {Eq. (12), p. 208}) \\ f \in \mathrm{L} ^ {2} ((\mathbf {R} / a \mathbf {Z}) ^ {n}) & \frac {1}{a ^ {n}} \int_ {(\mathbf {R} / a \mathbf {Z}) ^ {n}} | f (x) | ^ {2} d x = \sum_ {y \in \mathbf {Z} ^ {n}} | \widehat {f} (y) | ^ {2} \\ & (\text {Th. 1, p. 215}) \\ f \in \mathrm{L} ^ {2} ((\mathbf {R} / a \mathbf {Z}) ^ {n}) \text {such that} & f (x) = \sum_ {y \in \mathbf {Z} ^ {n}} \widehat {f} (y) \exp (\frac {2 i \pi}{a} y \cdot x) \\ \widehat {f} \in \mathrm{L} ^ {1} (\mathbf {Z} ^ {n}) & (\text {almost everywhere; Prop. 12, p. 219}) \\ f \in \mathrm{A} ((\mathbf {R} / a \mathbf {Z}) ^ {n}) & f (x) = \sum_ {y \in \mathbf {Z} ^ {n}} \widehat {f} (y) \exp (\frac {2 i \pi}{a} y \cdot x) \\ & (x \in (\mathbf {R} / a \mathbf {Z}) ^ {n}; \text {Prop. 11 p. 217}) \end{array}
\]

\specialsection{\texorpdfstring{Poisson's formula in \(\mathbf{R}^n\)}{Poisson's formula in R\^{}n}}

We state here the Poisson formula in the particularly important case where n = 1 (corollary of Proposition 15 of II, p.~230).

Let \(f\in \mathrm{L}^1 (\mathbf{R})\) and \(a > 0\) . Suppose that

\[
\sum_ {h \in \mathbf {Z}} | \widehat {f} (a ^ {- 1} h) | <   + \infty ,
\]

\[
\sum_ {h \in \mathbf {Z}} | f (x + a h) | <   + \infty \text {   for all   } x \in \mathbf {R},
\]

\[
x \mapsto \sum_ {h \in \mathbf {Z}} f (x + a h) \text {   continuous   on   } \mathbf {R}.
\]

Then, for all \(x \in \mathbf{R}\) , one has

\[
\sum_ {h \in {\bf Z}} f (x + a h) = \frac {1}{a} \sum_ {k \in {\bf Z}} \widehat {f} \Bigl (\frac {k}{a} \Bigr) \exp (\frac {2 i \pi}{a} k x).
\]

\specialchapter{INDEX OF NOTATIONS}

SpA(x) (spectrum) 2 R(x,λ) (resolvent) 3 Ā (adjoining of an identity element) 4 Sp′A(x) (spectrum) 4 X(A) (characters) 6 X(h) (character functor) 6 G\_A(x) (Gelfand transform) 7 X′(A) (characters) 9 X′(h) (character functor) 9 G′A(x) (Gelfand transform) 10 J(A) (primitive ideals) 12 V(M) (closed set of J(A)) 12 Υ(T) (ideal associated with T) 13 cl(π) (class of π) 13 Ā (set of irreducible representations) 14 γ (regular representation) 16 δ (regular representation) 16 B(X;K) (bounded functions) 17 C\_b(X;K) (bounded continuous functions) 17 C\_0(X;K) (continuous functions vanishing at infinity) 17 C(X;K) (continuous functions) 17 K(X;K) (continuous functions with compact support) 17 M¹(G) (bounded measures) 19 ρ(x) (spectral radius) 21 ev\_x (evaluation at x) 32 X′ (Alexandroff compactification) 33 Sp\^{}Λ\_A(x) (simultaneous spectrum) 42 O(U;F) (holomorphic functions) 49 O(K;F) (germs) 50 O(U;A) (algebra of holomorphic functions) 50



\(\mathcal{O}(\mathrm{K};\mathrm{A})\) (algebra of germs) 50 \(\Theta_{a}\) (holomorphic functional calculus) 51 \(f(a)\) (functional notation for the holomorphic functional calculus) 72 \(j_{\mathrm{H}}\) (idempotent associated with H) 82 \(\mathrm{R_H}(x,\lambda)\) (resolvent with respect to \(\mathrm{A_H} = j_{\mathrm{H}}\mathrm{A}j_{\mathrm{H}}\) ) 83 \(\mathcal{O}_{\mathbf{R}}(\mathrm{S})\) ( \(\mathbf{R}\) -algebra of germs \(f = f^{*}\) ) 86 \(\mathrm{Sp}_{\mathrm{A(C)}}(x)\) (complex spectrum) 86 \(x^{*}\) (adjoint) 96 \(\mathrm{A}_h\) (Hermitian elements) 96 \(\mathcal{C}'(\mathrm{X})\) (functions vanishing at a point) 103 \(f(x)\) (continuous functional calculus) 110 \(\mathrm{A}_{+}\) (positive elements) 115 \(x \geqslant y\) (order relation on \(\mathrm{A}_h\) ) 117 \(x^{+}\) (positive part) 117 \(x^{-}\) (negative part) 117 \(|x|\) (absolute value) 117 \(x^{\alpha}\) (power of a positive element) 118 \(\sqrt{x}\) (square root of a positive element) 118\\
Stell(A) (enveloping star algebra) 124\\
Stell(G) (star algebra of a group) 125 \(e_{\mathrm{H}}(u)\) (spectral projection) 129 \(\mathrm{E_H}(u)\) (spectral subspace) 129 \(\widetilde{\mathrm{E}}_{\mathrm{H}}(u)\) (spectral subspace) 129 \(\iota(u)\) (numerical range) 135 \(m(u)\) (lower bound of the spectrum) 138 \(\mathrm{M}(u)\) (upper bound of the spectrum) 138 \(|u|\) (absolute value of a morphism between Hilbert spaces) 139 \((j,|u|)\) (polar decomposition) 140\\
P(X) (limits of polynomials) 146 \(\operatorname{lsc}(g)\) (length of commutators) 155 \(\mathbf{C}[\mathrm{G}]\) (group algebra) 183 \(\operatorname{Stell}_r(\mathrm{G})\) (reduced stellar algebra) 184 \(\operatorname{Lip}_b(\mathrm{X})\) (bounded Lipschitz functions) 196 \(\widetilde{f}(x)\) (involution on \(\mathrm{L}^1 (\mathrm{G})\) ) 199 \(\widehat{\mathbf{G}}\) (dual group) 201 \(\chi (\mu)\) (character of \(\mathcal{M}^1 (\mathrm{G})\) ) 202 \(\mathrm{A}^{\perp}\) (orthogonal) 205 \(\widehat{\varphi}\) (dual morphism) 205 \(\mathcal{F}_{\mathrm{G}}(\mu)\) (Fourier transform) 207 \(\overline{\mathcal{F}}_{\mathrm{G}}(\mu)\) (inverse Fourier transform) 207 \(\mathcal{F}_{\mathrm{G}}(f)\) (Fourier transform) 208 \(\overline{\mathcal{F}}_{\mathrm{G}}(f)\) (inverse Fourier transform) 208\\
A(G) (subspace generated by \(f*g\) ) 210 \(\widehat{\mathrm{A}}(\mathrm{G})\) (image of A(G) under \(\mathcal{F}\) ) 211

B(G) (subalgebra of \(f \in L^{1}(G)\) with \(\mathcal{F}(f) \in L^{1}(\widehat{G})\) ) 222 T (torus) 235 \(\mathcal{S}(\mathbf{R}^{n})\) (Schwartz functions) 238 V(Λ) (covolume) 239 \(\mathcal{S}'(\mathbf{R}^{n})\) (tempered distributions) 239 A(W) (characters in W) 258 Y(F) (subspace of L∞(G) generated) 258 s(t) (sign function) 261 \(\mathcal{P}(X)\) (probability measures) 261 ζ(s) (zeta function) 267 L(s, χ) (Dirichlet L-function) 268 L(G) (representations of G in R) 307 Λ(n) (von Mangoldt function) 312

\specialchapter{TERMINOLOGICAL INDEX}

\textbf{A}

Adjoint, 96 Adjoining of an identity element, 4, 105 Banach algebra, 15 involutive Banach algebra, 100 regular Banach algebra, 89 involutive, 96 normed, 15 normed involutive, 100 normed product, 16 involutive subalgebra, 97 unital, 1 Star algebra, 102 of a group, 125 enveloping, 124 quotient, 122 reduced, 184 μ-ergodic map, 190 Proper partial map, 33 Self-adjoint (--- part), 96

\textbf{B}

Bernoulli (--- shift), 191 Bernstein (--- inequality), 279 Bohr-Pál (--- theorem), 286 Shilov boundary, 171

\textbf{C}

Continuous functional calculus, 110 holomorphic, 51

holomorphic in one variable, 74 Character, 6, 9 Dirichlet, 268 Legendre, 288 unital, 6 unitary, 201 Cauchy-Davenport (--- theorem), 276 Chebotarev (--- theorem), 276 Cohen factorization theorem, 185 idempotent theorem, 272 Ditkin condition, 92 spectral synthesis condition, 316 Symmetric convergence, 261 Convex (polynomially --- set), 44 Fourier cotransformation, 207 Fourier cotransform, 207 Covolume, 239 Weyl's equidistribution criterion, 277

\textbf{D}

Bernoulli shift, 191 Polar decomposition, 140 in a star algebra, 119 Spectral decomposition, 130 Density of a sphere packing,

Dirichlet character ---, 268 --- L-function, 268 --- theorem, 269 Tempered distributions, 239 Topological divisor of zero, 23 Dual group ---, 201 --- measure, 214 --- morphism, 205 --- lattice, 238 Pontryagin duality, 220 between groups, 225

\textbf{E}

Positive element, 115 Sphere packing, 295 Resolvent set, 2 Polynomially convex hull, 45 Envelopment, 70 Schrödinger equation, 269 Equidistribution modulo 1, 277 Ergodic (μ--- map), 190 Space of characters, 8 Exponential, 78

\textbf{F}

Fejér (--- theorem), 242 Fekete (--- lemma), 20 Fermat (--- theorem), 288 Faithful (--- representation), 11 Characteristic function of a probability measure, 280 conjugate-harmonic, 280 Bessel function of the first kind, 269 Möbius function, 313 Schwartz function, 238 von Mangoldt function, 312 Hermite function, 263 Dirichlet L-function, 268 slowly oscillating, 253 Lipschitz, 196 almost periodic, 292 sign, 261 theta, 297 Riemann zeta, 267

Associated differential form, 52 Hermitian linear form, 98 positive, 183 Modular form, 299 Hasse--Davenport formula, 288 Plancherel formula, 217 Poisson formula, 231 Fourier inversion formula, 219 Fourier cotransformation of ---, 207 --- inversion formula, 219 --- series, 241 --- transformation, 207 Fuglede--Putnam (theorem), 106

\textbf{G}

Gauss (sums of ---), 287, 299 Gaussian (law ---), 281 Gelfand (transformation of ---), 7, 10 Gelfand--Mazur (theorem of ---), 26 Germs of holomorphic functions, 49 Gibbs (phenomenon of ---), 299 Dual group, 201 monothetic, 306 amenable, 291 solenoidal, 306 Groups in duality, 225

\textbf{H}

Hadamard--de la Vallée Poussin (theorem of ---), 313 Hardy (uncertainty principle), 274 Hardy--Littlewood (theorem of ---), 309 Hasse--Davenport (formula of ---), 288 Hausdorff--Toeplitz (theorem of ---), 136 Hermitian element ---, 96 linear form ---, 98 Hildebrandt (theorem of ---), 189

\textbf{I}

Primitive ideal, 11 regular, 10

Idempotent associated with an element, 82 Ikehara (tauberian theorem of), 312 Numerical image of an endomorphism, 135 Inequality of Bernstein, 279 of large sieve, 294 of van der Corput, 277 isoperimetric, 284 Involution, 96 semilinear, 95 Involutive algebra ---, 96 Banach algebra ---, 100 normed algebra ---, 100 Irreducible (representation ---), 11

\textbf{J}

Jacobi (sums of ---), 288 symbol of ---, 299 Jacobson (topology of ---), 13

\textbf{K}

Kaplansky (theorem of ---), 184 Kolmogorov (theorem of ---), 290 Kronecker (theorem of ---), 277

\textbf{L}

Laplace (transformation of ---), 312 Legendre (character of ---), 288 Lemma of Fekete, 20 Lévy theorem of ---, 67 continuity theorem of ---, 281 Logarithm, 78 Poisson law, 282 of quadratic reciprocity, 297 Gaussian, 281 Stable length of commutators,

\textbf{M}

Counting measure, 200 normalized Haar measure, 200 probability measure, 261 dual measure, 214 Morphism

of involutive algebras, 97 of star algebras, 102 of representations, 11 unital, 1

\textbf{N}

Pisot numbers, 295 Prime numbers (theorem of ---), 313 Normal (element), 96

\textbf{O}

Oka--Weil (theorem of ---), 68 Orthogonal, 205 Orthogonal projector, 133

\textbf{P}

Paley-Wiener (theorem of ---), 278\\
Permanent, 156\\
Phenomenon of Gibbs, 299\\
Plancherel\\
formula of ---, 217\\
theorem of ---, 215\\
Full (subalgebra ---), 5\\
Poisson\\
formula of ---, 231\\
law of ---, 282\\
Polynomially convex\\
set ---, 44\\
envelope ---, 45\\
Pontryagin (duality theorem of ---), 220\\
Positive\\
element ---, 115\\
systematically --- element, 182\\
Primitive (ideal ---), 11\\
Principle\\
of uncertainty, 273, 275\\
of uncertainty of Hardy, 274\\
Eulerian product, 267\\
Arithmetic progression (theorem of the --- of Dirichlet), 269\\
Arithmetic progressions of length 3, 271\\
Spectral projector, 129\\
Proper (partial map ---), 33

\textbf{Q}

Quasi-nilpotent (element ---), 21

R Rademacher-Menshov (theorem of ---), 287 Spectral radius, 21 of a subset, 157 Regular (ideal ---), 10 Regular Banach algebra ---, 89 representation ---, 16 Orthogonality relations, 233 Representation of an algebra, 11 faithful, 11 irreducible, 11 regular, 16 Equivalent representations, 11 Lattice covolume of a ---, 239 dual, 238 Resolvent, 3 pole of the ---, 131 Riemann (zeta function of ---), 267 Roth (theorem of ---), 271 Runge (theorem of ---), 69, 151

\textbf{S}

Sahakyan (theorem of ---), 285 Salem (theorem of ---), 295 Schwartz (functions of ---), 238 Semilinear (involution ---), 95 Fourier series, 241 Gauss sums, 287, 299 Jacobi sums, 288 Salié sums, 288 Involutive subalgebra, 97 full, 5 Independent subset, 265 Spectral subspace, 129 One-parameter subgroup, 307 Spectral decomposition ---, 130 projector ---, 129 radius ---, 21 subspace ---, 129 Spectrum, 2, 4 complex, 86 of an endomorphism, 127

exponential, 173 simultaneous, 42 singular, 158 Adapted sequence, 52 of Følner, 291 Exact sequence of topological groups, 225 dual, 226 Jacobi symbol, 299 Topological generating system, 16

\textbf{T}

Theorem of Bohr-Pál, 286 of Cauchy-Davenport, 276 of Chebotarev, 276 of continuity of Lévy, 281 of duality of Pontryagin, 220 of factorization of Cohen, 185 of Fejér, 242 of Fermat, 288 of Fuglede--Putnam, 106 of Gelfand--Mazur, 26 of Hausdorff--Toeplitz, 136 of Hildebrandt, 189 of Kaplansky, 184 of Kolmogorov, 290 of Kronecker, 277 of the arithmetic progression of Dirichlet, 269 of Oka--Weil, 68 of P. Lévy, 67 of Paley-Wiener, 278 of Plancherel, 215 of Rademacher-Menshov, 287 of Roth, 271 of Runge, 69, 151 of Sahakyan, 285 of Salem, 295 of Wiener, 38 of idempotents of Cohen, 272 of prime numbers, 313 ergodic of von Neumann, 190 fundamental limit of probability theory, 281 tauberian of Hardy--Littlewood, 309 tauberian of Ikehara, 312

tauberian of Wiener, 253 Topology of Jacobson, 13 weak, 8, 10 Transformation of Fourier, 207 of Gelfand, 7 of Hankel, 270 of Laplace, 312 Fourier transform, 207

\textbf{U}

Unital algebra ---, 1 morphism ---, 1 Unitary

character ---, 201 element ---, 96 Approximate unit, 120 increasing, 120

V Eigenvalue, 128 van der Corput (inequality of ---), 277 von Neumann (ergodic theorem of ---), 190

W Weyl (criterion of ---), 277 Wiener theorem of ---, 38 tauberian theorem of ---, 253
